\documentclass[11pt,a4paper]{article}
\usepackage[T1]{fontenc}
\usepackage{lmodern}
\usepackage[margin=25mm,headheight=14pt,headsep=7mm,footskip=11mm]{geometry}
\usepackage{amsmath,amssymb,booktabs,array,tabularx,microtype,xcolor,fancyhdr}
\usepackage{enumitem,listings,needspace,tocloft}
\usepackage{mathpartir,stmaryrd,tcolorbox,caption,placeins}
\usepackage[hidelinks]{hyperref}
\hypersetup{pdftitle={Learning to make the proof explicit},
  pdfauthor={Oliver Bracevac},
  pdfsubject={A worked tutorial on dependent object types, explicit evidence, and captures}}
\definecolor{ink}{RGB}{29,48,68}
\definecolor{wash}{RGB}{239,244,248}
\definecolor{hint}{RGB}{74,95,108}
\definecolor{srccol}{RGB}{0,70,150}
\definecolor{tgtcol}{RGB}{160,40,30}
\definecolor{runcol}{RGB}{45,48,52}
\newcommand{\src}[1]{{\color{srccol}#1}}
\newcommand{\tgt}[1]{{\color{tgtcol}#1}}
\newcommand{\run}[1]{{\color{runcol}#1}}
\newcommand{\tr}[1]{\llbracket\src{#1}\rrbracket}
\newtcolorbox{examplepanel}[2]{colback=#1!3!white,colframe=#1!50!white,
  colbacktitle=#1!8!white,coltitle=#1,title={#2},
  fonttitle=\small\sffamily\bfseries,boxrule=0.4pt,arc=1mm,
  left=8pt,right=8pt,top=5pt,bottom=5pt,before skip=5pt,after skip=5pt}
\captionsetup[figure]{font=small,labelfont={bf,sf},labelsep=period,
  justification=raggedright,singlelinecheck=false,skip=7pt}
\setlength{\textfloatsep}{16pt plus 3pt minus 2pt}
\setlength{\intextsep}{14pt plus 3pt minus 2pt}
\renewcommand{\topfraction}{0.92}
\renewcommand{\bottomfraction}{0.85}
\renewcommand{\textfraction}{0.08}
\renewcommand{\floatpagefraction}{0.75}
\pagestyle{fancy}
\fancyhf{}
\fancyhead[L]{\small\sffamily\color{ink} LEARNING TO MAKE THE PROOF EXPLICIT}
\fancyhead[R]{\small\sffamily A worked tutorial}
\fancyfoot[L]{\small\sffamily Oliver Bra\v{c}evac}
\fancyfoot[R]{\small\sffamily\thepage}
\renewcommand{\headrulewidth}{0.3pt}
\setlength{\parindent}{0pt}
\setlength{\parskip}{6pt plus 1pt minus 1pt}
\setlength{\emergencystretch}{2em}
\setlist{nosep,topsep=5pt,leftmargin=1.5em}
\raggedbottom
\lstset{basicstyle=\ttfamily\small,columns=fullflexible,keepspaces=true,
  frame=none,breaklines=true,xleftmargin=1em,aboveskip=8pt,belowskip=8pt}
\lstdefinestyle{source}{basicstyle=\ttfamily\small\color{srccol}}
\lstdefinestyle{target}{basicstyle=\ttfamily\small\color{tgtcol}}
\lstdefinestyle{runtime}{basicstyle=\ttfamily\small\color{runcol}}
\newcommand{\FC}{\ensuremath{\mathrm{FC}_{\mathrm{dot}}}}
\newcommand{\ev}[1]{\tgt{\mathsf{#1}}}
\newcommand{\ty}[1]{\mathit{#1}}
\newcommand{\cast}[2]{\ev{cast}(#1,#2)}
\newcommand{\erase}[1]{\lfloor #1\rfloor}
\newcommand{\obj}[1]{\mu s.[\,#1\,]}
\newcommand{\iface}[2]{\src{\{A:#1\,..\,#2\}}}
\newcommand{\roots}[1]{\operatorname{roots}_{\sigma}(#1)}
\newcommand{\uses}[1]{\operatorname{uses}(#1)}
\newcommand{\view}[1]{\operatorname{view}_{\sigma}(#1)}
\newcommand{\hnf}[1]{\operatorname{hnf}_{\sigma}(#1)}
\newcommand{\form}[1]{\tgt{\widehat{\mathsf{#1}}}}
\newcommand{\lesson}[1]{\FloatBarrier\par\addvspace{10pt}\needspace{10\baselineskip}\section{#1}}
\newcommand{\focus}[1]{\noindent{\itshape\color{hint}#1}\par\medskip}
\newcommand{\remember}[1]{\par\smallskip\noindent
  \colorbox{wash}{\parbox{\dimexpr\linewidth-2\fboxsep\relax}{\strut#1\strut}}
  \par\smallskip}
\newcounter{exercise}
\newcommand{\exercise}[2]{\refstepcounter{exercise}\label{ex:#1}
  \par\needspace{4\baselineskip}\textbf{Exercise \theexercise.} #2
  \hfill\mbox{\small\hyperref[ans:#1]{Answer on p.~\pageref*{ans:#1}}}\par}
\newcommand{\answer}[1]{\par\needspace{4\baselineskip}
  \phantomsection\label{ans:#1}\textbf{Exercise \ref{ex:#1}.}\ }
\newcommand{\leanfile}[2]{\href{https://github.com/bracevac/path-dependent/blob/main/lean/Coercions/#1}{\texttt{#2}}}

\begin{document}
\thispagestyle{fancy}
{\huge\sffamily\bfseries\color{ink} Learning to make\\[3pt] the proof explicit\par}
\vspace{6pt}
{\Large Dependent object types, executable evidence, and captures\par}
\vspace{10pt}

Dependent object types let a client use a member type $x.A$ while
knowing only its bounds. Proving type safety requires connecting those
interface assumptions to the definitions and values that an execution
can supply.

This tutorial develops a proof architecture in which source typing
derivations translate to target terms, with subtyping derivations
represented by explicit evidence. Target typing checks that evidence;
normalization interprets
it against a typed store and exposes the structure needed for execution.
The running examples follow this route from source derivations through
coercions, substitution, and normalization to individual machine steps.
The capture extension applies the same method to subcapturing and
effect safety.

\textbf{Prerequisites.}
We assume familiarity with typed lambda calculi, typing and subtyping
derivations, substitution, and small-step semantics. The tutorial
introduces the specific features of DOT, \FC{}, and capture checking,
and explains the Lean notation used in its executable companion.

\textbf{The worked examples.}
The core example constructs an object-interface coercion, tracks the
dependent name it preserves, and uses its normalized evidence to execute
a function call. The capture examples distinguish a computation's use
set from the capture annotation of its result and derive the root
inclusions needed for safety.

\textbf{How to read.}
Sections~\ref{sec:objects}--\ref{sec:bad} follow one object and one function.
Section~\ref{sec:proof} connects the calculations to safety.
Sections~\ref{sec:capture}--\ref{sec:capture-scope} repeat the argument
for captures. Section~\ref{sec:larger} then explains which further
object features require additional cases. The final sections connect the account to
the implementation and give worked solutions. The exercises isolate
individual proof obligations and calculations from the examples.

\remember{\textbf{The central move.}
Schematically, a derivation of $S<:T$ translates to evidence $e:S\le T$.
Subtyping becomes a target proposition, and its derivation becomes an
explicit inhabitant. Checking this evidence under assumptions and
interpreting it over stored values are different obligations. Substitution
connects them: it replaces the assumptions by actual typed atoms. This
is the bridge we will use both for type safety and for capture prediction.}

\textbf{Relation to the development.}
All metatheoretic claims refer to the existing Lean development and its
two reference reports. We present selected constructs in named notation,
then relate that notation to the implementation.
Section~\ref{sec:lean} identifies the executable companion, the checked
examples, and the scope of the expository simplifications.

\clearpage
\setcounter{tocdepth}{1}
\setlength{\cftbeforesecskip}{3pt}
\renewcommand{\cftsecfont}{\small}
\renewcommand{\cftsecpagefont}{\small}
\begingroup
\setlength{\parskip}{0pt}
\tableofcontents
\endgroup
\thispagestyle{fancy}
\bigskip
\textbf{Notation reference.}
The text introduces each notation before its use in a calculation.
\begin{center}
\renewcommand{\arraystretch}{1.25}
\begin{tabularx}{\linewidth}{@{}lX@{}}
\toprule
Notation & Read it as\\
\midrule
$S<:T$ & A source subtyping judgment.\\
$e:S\le T$ & Target evidence $e$ for an inclusion.\\
$q:S\equiv T$ & Target evidence $q$ for equality.\\
$p;q$ & First use $p$, then use $q$; their middle endpoints must agree.\\
$x.A$ & The stable name of type member $A$ belonging to $x$.\\
$\cast{a}{e}$ & The same underlying atom presented using evidence $e$.\\
$\obj{P_0,\ldots,P_n}$ & An object interface listing facts about its self $s$.\\
$\sigma$ & The store of allocated values.\\
$\form{id}$ & A head form whose endpoints agree after resolving definitions.\\
$C\le_c D$ & A subcapturing judgment, also represented by evidence.\\
\bottomrule
\end{tabularx}
\end{center}
\small
\textbf{Color key.} \src{Blue: source terms and derivations.}
\tgt{Red: target terms and evidence.}
\run{Charcoal: store facts and erased computation.}
Figure panels also name their role, so the distinction survives grayscale printing.
\normalsize

\clearpage\lesson{Type members and abstraction}\label{sec:objects}
\focus{A type member separates an object's type definition from the bounds available to its clients.}

In Dependent Object Types (DOT), the source calculus family studied here,
an object can define a type member as well as hold value fields.
For a type $W$, consider the source-style sketch
\begin{lstlisting}
let x = { type A = W }
\end{lstlisting}
The selection $x.A$ names the type member of this particular object.
Another object $y$ may define a member with the same label $A$, but the
selections $x.A$ and $y.A$ are distinct names. Relating them requires
information about their definitions or bounds.

For the running example, fix
\[
  W=\top\to\top,\qquad f=\lambda(u:\top).u.
\]
Here $\top$ is the universal type, so $f:W$. We also use $\bot$ for the
empty type: a closed well-typed computation cannot return a value of
that type, although it may diverge. Using a function type for $W$ keeps
the example within the formal calculus, without additional primitives.

\subsection*{Definitions and interfaces}
The declaration $A=W$ defines the member $A$ as $W$. An interface
$\iface{S}{T}$ specifies the bounds available to a client:
\[
  x:\iface{S}{T}
  \quad\Longrightarrow\quad
  S<:x.A\quad\hbox{and}\quad x.A<:T.
\]
The lower bound permits a value of type $S$ to be used at the abstract
type $x.A$; the upper bound permits a value of type $x.A$ to be used
at $T$.

The exact interface $\iface{W}{W}$ provides both $W<:x.A$ and $x.A<:W$.
The less informative interface $\iface{\bot}{\top}$ provides only
$\bot<:x.A$ and $x.A<:\top$. The object can have the same definition in
both cases; widening the interface restricts the facts available to the
client.

This is the abstraction provided by a module that exposes a type and
operations on it while hiding its representation. The running object
contains only a type definition, allowing us to study the reasoning
about type members before introducing fields and methods.

\subsection*{The running program}
Figure~\ref{fig:program} gives the running program and its three subtyping
obligations. Subsequent sections translate its derivation, construct the
corresponding store, and execute the final application.


\begin{figure}[!htbp]
\begin{examplepanel}{srccol}{Source program \(P\) --- the running example}
\src{\(W=\top\to\top\). Let annotations specify the assigned types.}
\begin{lstlisting}[style=source]
let x         = { type A = W }
let f         = lambda (u: Top). u
let hidden    : x.A = f          -- L: W <: x.A
let recovered : W   = hidden     -- H: x.A <: W
recovered f                     -- T: W <: Top
\end{lstlisting}
\end{examplepanel}
\begin{examplepanel}{srccol}{Subtyping obligations}
\[
\src{
 \begin{array}{ccl}
 \text{L: hide} & W<:x.A & \text{the lower bound of }x.A,\\
 \text{H: recover} & x.A<:W & \text{the upper bound of }x.A,\\
 \text{T: pass the argument} & W<:\top & \text{inclusion into the universal type}.
 \end{array}}
\]
\end{examplepanel}
\caption{The running source program, with the object viewed at
\(\iface{W}{W}\). The aliases present the identity first at \(x.A\)
and then at \(W\); the final application presents its argument at
\(\top\). This sketch exposes the member bounds directly.
Section~\ref{sec:evidence} translates its client body;
Section~\ref{sec:closed-program} gives the complete translation, including
the recursive opening required by the formal source.}
\label{fig:program}
\end{figure}



The binding of $\ty{hidden}$ uses the lower bound $W<:x.A$.
The binding of $\ty{recovered}$ uses the upper bound $x.A<:W$.
The final application uses $W<:\top$ to type its argument. Since both
$\ty{recovered}$ and $f$ denote the identity, evaluation returns the
identity function supplied as the argument.

The displayed program is an explanatory sketch, rather than accepted
Scala syntax. The formal source names intermediate expressions and
explicitly records the step that exposes an object literal's member
interface; these details are introduced where the translation needs them.

\remember{The stored definition $A=W$, the stable name $x.A$, and the
interface supplying facts about that name have distinct roles. An
interface change preserves the definition and the name.}

\exercise{bounds}{At the abstract interface $\iface{\bot}{\top}$, can a
client justify assigning our $f:W$ to a variable of type $x.A$ using the
bounds alone? Write the inclusion it would need and compare it with
the available lower bound.}

\lesson{From typing derivations to target terms}\label{sec:evidence}
\focus{The translation represents subtyping derivations as evidence and subsumption as a cast.}

The source subsumption rule combines a typing derivation $h$ with a
subtyping derivation $d$:
\[
\src{\frac{h:\Gamma\vdash u:S\qquad d:\Gamma\vdash S<:T}
           {\textsc{Sub}(h,d):\Gamma\vdash u:T}}.
\]
Subsumption changes the assigned type while preserving the source
expression $u$. The chosen conversion is therefore recorded in the
derivation rather than the term.

Write $\tr{h}$ for the target term translated from a source typing
derivation $h$, and $\tr{d}$ for the evidence translated from a source
subtyping derivation $d$. The subsumption clause is
\[
\tgt{\tr{\textsc{Sub}(h,d)}=\cast{\tr{h}}{\tr{d}}.}
\]
This translation clause recursively translates the two premises and
combines their results with the target $\ev{cast}$ constructor.

The target is called \FC{}. An expression $e:S\le T$ is called
\emph{inclusion evidence}, or an \emph{inclusion coercion}.
The target has no implicit subsumption: a cast records a type change
justified by evidence, with no potentially failing runtime type test.
Its typing rule checks the two translated premises:
\[
\tgt{\frac{\Gamma'\vdash \tr{h}:S'\qquad
           \Gamma'\vdash\tr{d}:S'\le T'}
          {\Gamma'\vdash\cast{\tr{h}}{\tr{d}}:T'}}.
\]
Primes denote the translated context and types. We use the same notation
for $W$ and $x.A$ on both sides of the translation.

A target \emph{atom} is a variable, possibly surrounded by typing
wrappers such as casts. Its \emph{root} is the underlying variable:
both $f$ and $\cast{f}{e}$ have root $f$. Atoms serve as the receivers
and arguments of operations.

The \emph{erasure} $\erase{t}$ removes typing annotations and evidence
wrappers, yielding the underlying program. In particular,
$\erase{\cast{u}{e}}=\erase{u}$.
Erasure relates program syntax; the typed target's operational semantics
may still consult evidence.

\subsection*{Translating the two bound selections}
Consider an occurrence of $f$ under the parameter context
$\src{\Gamma=x:\iface{W}{W},\,f:W}$.
This is the client part of the example: for now, assume that its object
parameter already exposes the exact bounds. We will translate this body
first. Section~\ref{sec:closed-program} then supplies the literal's
introduction and recursive opening, displays the complete target program,
and substitutes the allocated atoms into this same body. No assumption
about a stored definition is needed for the present calculation.
Reading the assumed interface gives two source derivations:
\[
\src{\begin{aligned}
d_\ell&:\Gamma\vdash W<:x.A,\\
d_h&:\Gamma\vdash x.A<:W.
\end{aligned}}
\]
Write the translated evidence using the abbreviations
\[
 \ell_x=\ev{lower}_A(x):W\le x.A,\qquad
 h_x=\ev{upper}_A(x):x.A\le W.
\]
Section~\ref{sec:telescopes} expands each abbreviation into an indexed
member read. Figure~\ref{fig:derivation} translates the two subsumption
steps from their premises: variable typing yields the atom $f$ and the
lower-bound derivation yields $\ell_x$, which the inner \textsc{Sub}
combines into $\cast{f}{\ell_x}$. The upper-bound derivation then yields
$h_x$, and the outer \textsc{Sub} adds the second cast.


\begin{figure}[!htbp]
\begin{examplepanel}{srccol}{Source derivation --- two subsumptions on the same occurrence of \(f\)}
\small
\src{\(\Gamma=x:\iface{W}{W},\,f:W\).
Here \(h_f\) derives variable typing for \(f\), and \(d_\ell,d_h\)
derive the lower and upper bounds of \(x.A\).}
\[
\src{
\inferrule*[right=\textsc{Sub}]{
 \inferrule*[right=\textsc{Sub}]{
 h_f:\Gamma\vdash f:W\\
 d_\ell:\Gamma\vdash W<:x.A
 }{\Gamma\vdash f:x.A}\\
 d_h:\Gamma\vdash x.A<:W
}{\Gamma\vdash f:W}}
\]
\end{examplepanel}
\begin{examplepanel}{tgtcol}{Translation of the derivation}
\[
\tgt{
\begin{aligned}
 \tr{h_f}&=f,\\
 \tr{d_\ell}&=\ell_x=\ev{lower}_A(x):W\le x.A,\\
 \tr{d_h}&=h_x=\ev{upper}_A(x):x.A\le W,\\
 \tr{\textsc{Sub}(h_f,d_\ell)}&=\cast{f}{\ell_x},\\
 \tr{\textsc{Sub}(\textsc{Sub}(h_f,d_\ell),d_h)}
   &=\cast{\cast{f}{\ell_x}}{h_x}.
\end{aligned}}
\]
\end{examplepanel}
\begin{examplepanel}{runcol}{Erasure of the target term}
\[
\run{\erase{\cast{\cast{f}{\ell_x}}{h_x}}=f.}
\]
\end{examplepanel}
\caption{Each \textsc{Sub} node contributes a target \(\ev{cast}\),
while the source expression remains \(f\) throughout the derivation.
The calculation assumes parameter binders for \(x\) and \(f\).
Figure~\ref{fig:let-translation} extends it to the program's let bindings.}
\label{fig:derivation}
\end{figure}



\subsection*{Dependence on the source derivation}
The judgment $\src{\Gamma\vdash f:W}$ admits several derivations.
Variable typing translates to $f$, while the two subsumptions translate
to $\cast{\cast{f}{\ell_x}}{h_x}$. A third derivation first composes
the bounds:
\[
\src{\frac{d_\ell:\Gamma\vdash W<:x.A\qquad
           d_h:\Gamma\vdash x.A<:W}
          {\textsc{Trans}(d_\ell,d_h):\Gamma\vdash W<:W}}.
\]
Source transitivity becomes an evidence constructor:
\[
\tgt{\tr{\textsc{Trans}(d_\ell,d_h)}
       =\ev{trans}(\ell_x,h_x).}
\]
One subsequent subsumption therefore gives
$\cast{f}{\ev{trans}(\ell_x,h_x)}$.
The three target expressions are syntactically distinct, although all
erase to $f$. In particular, nested casts and a cast by composed evidence
are different constructor expressions. The translation is thus a
function of the source derivation, not only of its concluding term.

We abbreviate $\ev{trans}(p,q)$ by $p;q$, in derivation order:
\[
 \frac{p:S\le M\qquad q:M\le T}{p;q:S\le T}.
\]
Composition requires matching intermediate endpoints. Here the chain is
$W\xrightarrow{\ell_x}x.A\xrightarrow{h_x}W$.

\subsection*{Let translation and evaluation}
The translation preserves the running program's let bindings:
\[
\tgt{
\tr{\textsc{Let}(h_1,h_2)}
=\ev{let}\ y=\tr{h_1}\ \ev{in}\ \tr{h_2}.}
\]
Figure~\ref{fig:let-translation} applies this clause to the body of the
running example. The first right-hand side uses the lower-bound cast.
The second uses the upper-bound cast on the newly bound name
$\ty{hidden}$. The application also needs the argument cast
$\ev{top}(W):W\le\top$.

Application translates the variable typing derivations of its operands
to atoms. Subsumption in such a derivation contributes a cast wrapper,
so the source argument $f$ becomes $\cast{f}{\ev{top}(W)}$.
The full language distinguishes casts of terms, values, and atoms;
we use the same notation when the syntactic position determines the sort.


\begin{figure}[!htbp]
\begin{examplepanel}{srccol}{Source body --- after binding \(x\) and \(f\)}
\begin{lstlisting}[style=source]
let hidden    : x.A = f in
let recovered : W   = hidden in
recovered f
\end{lstlisting}
\end{examplepanel}
\begin{examplepanel}{tgtcol}{Target translation with explicit casts}
\tgt{\(\ell_x=\ev{lower}_A(x)\), \(h_x=\ev{upper}_A(x)\),
\(t_W=\ev{top}(W):W\le\top\).}
\begin{lstlisting}[style=target,mathescape=true]
let hidden    = cast(f, $\ell_x$) in
let recovered = cast(hidden, $h_x$) in
recovered (cast(f, $t_W$))
\end{lstlisting}
\end{examplepanel}
\begin{examplepanel}{tgtcol}{Evaluation by substitution of the aliases}
\small
\[
\tgt{\begin{aligned}
 &\ev{let}\ \ty{recovered}=\cast{\cast{f}{\ell_x}}{h_x}\ \ev{in}\
       \ty{recovered}\ \cast{f}{t_W}\\
 &\hspace{5em}\longrightarrow^{*}
       \cast{\cast{f}{\ell_x}}{h_x}\ \cast{f}{t_W}.
\end{aligned}}
\]
The first binding substitutes \(\cast{f}{\ell_x}\) for \(\ty{hidden}\);
the second substitutes the nested cast for \(\ty{recovered}\).
The relation \(\longrightarrow^{*}\) includes the administrative steps
that rearrange pending casts.
\end{examplepanel}
\caption{Translation preserves the source let structure; evaluation
produces nested casts by substituting the let-bound atoms. The resulting
application erases to \(f\,f\). As in Figure~\ref{fig:derivation}, the
exact interface is assumed directly.
Section~\ref{sec:closed-program} adds the literal's introduction and
recursive opening, then instantiates its binders with atoms over the store.}
\label{fig:let-translation}
\end{figure}



During evaluation, substitution of the let-bound atoms produces the
nested wrappers. This display still has the parameters $x$ and $f$ free.
Closing them is a further calculation, given in
Section~\ref{sec:closed-program}; it must justify both the object
interface used by the evidence and the function supplied to the call.

\subsection*{Checking evidence endpoints}
The expression $\ell_x;\ev{top}(W)$ fails the composition rule:
its first component ends at $x.A$ and its second starts at $W$.
Inserting the upper-bound evidence gives the well-typed composition
$\ell_x;h_x;\ev{top}(W)$. Alternatively,
$\ell_x;\ev{top}(x.A)$ reaches $\top$ directly from $x.A$.

Even if a later store resolves $x.A$ to $W$, checking still compares
syntactic endpoints. Store-based normalization establishes what the
evidence means at runtime.

\subsection*{Erasure and operational correspondence}
The erasure equation $\erase{\cast{u}{e}}=\erase{u}$ helps show that
translation preserves the underlying erased program.
The typed target machine still consults a function atom's
evidence to perform a call. A simulation theorem relates those typed
steps to the erased program's execution.

\remember{Translation constructs evidence from source derivations;
evaluation substitutes atoms and consults their evidence; erasure
removes the evidence to relate the underlying programs.}

\exercise{order}{What does $h_x;\ell_x$ prove? Is it the same expression
as $\ell_x;h_x$? Also give explicit evidence for $W\le\top$ that goes
through $x.A$.}

\exercise{translate}{The source occurrence $f:W$ is first used at
$x.A$, then at $\top$ by the source \textsc{Top} rule $x.A<:\top$.
Give its target term when there are two \textsc{Sub} nodes, and when
there is one \textsc{Sub} over a \textsc{Trans} node. What does each erase to?}


\lesson{Interfaces as telescopes}\label{sec:telescopes}
\focus{The target represents an object interface as indexed propositions under a shared self binder.}

The interface $\iface{W}{W}$ specifies two bounds on its object's type
member. Its target representation binds a self variable $s$ in those
propositions:
\[
 \ty{Exact}=\obj{W\le s.A,\ s.A\le W}.
\]
This type describes an object satisfying the listed propositions, with
$s$ denoting the object itself. Here $\mu$ binds self in a proposition
list; this example does not involve recursive unfolding into a larger
type structure.

The development calls this list a \emph{telescope}. Its propositions are
indexed from zero:
\[
 \begin{array}{ccl}
 \text{position }0&:&W\le s.A,\\
 \text{position }1&:&s.A\le W.
 \end{array}
\]
All propositions share the self binder; individual propositions do not
bind term variables. Indexed selection therefore represents both lower
and upper bounds without a dedicated constructor for either.

\subsection*{Selecting a proposition}
For an atom $a:\ty{Exact}$, inclusion evidence is obtained with
\[
 \ev{member}\ a\ (\ev{refl}\ \ty{Exact})\ i.
\]
The atom $a$ identifies the object, the coercion selects an interface
reachable from its current type, and the index $i$ selects a proposition
in that interface. Since $a$ already has type $\ty{Exact}$, the coercion
is reflexivity, $\ev{refl}\ \ty{Exact}:\ty{Exact}\le\ty{Exact}$.

If $a$ is rooted at $x$, reading position 1 replaces $s$ by $x$:
\[
 \ev{member}\ a\ (\ev{refl}\ \ty{Exact})\ 1:x.A\le W.
\]
This expands the abbreviation $\ev{upper}_A(a)$. Position 0 similarly
gives $\ev{lower}_A(a):W\le x.A$. The abbreviations suppress only the
index and reflexivity evidence.

\subsection*{Roots and dependent substitution}
For the current fragment, atoms and their roots are defined by
\[
 a::=x\mid\cast{a}{e},\qquad
 \ev{root}(x)=x,\quad
 \ev{root}(\cast{a}{e})=\ev{root}(a).
\]
Suppose $a=\cast{x}{c}:\ty{Exact}$. The selected type is still $x.A$.
There is no new type name $(\cast{x}{c}).A$ in the type syntax.

When a function body checked under
$y:\ty{Exact}$ receives the atom $\cast{x}{c}$, occurrences of
$y$ as an atom become $\cast{x}{c}$, while occurrences of $y.A$ in
types become $x.A$. The wrapper supplies the required interface; the
root supplies the name.

Preserving the root ensures that the lower- and upper-bound proofs
continue to refer to the same selection after an interface change.
Their shared endpoint $x.A$ is what permits their composition in the
running example.

\subsection*{Equality propositions}
The literal that defines $A=W$ has the more precise type
\[
 \ty{Lit}=\obj{s.A\equiv W}.
\]
Its single equality proposition yields the two inclusions by reading
the equality in opposite directions. The fully abstract interface is
\[
 \ty{Abs}=\obj{\bot\le s.A,\ s.A\le\top}.
\]
The next two sections construct and evaluate the conversions
$\ty{Lit}\le\ty{Exact}\le\ty{Abs}$.

\exercise{root}{Let $a=\cast{\cast{x}{c_1}}{c_2}:\ty{Abs}$.
What are the endpoints obtained by reading positions 0 and 1?
Which part of the answer depends on the casts, and which part does not?}

\lesson{Constructing interface coercions}\label{sec:recipes}
\focus{An object coercion derives each destination proposition from the source telescope.}

The two interface conversions require evidence
\[
 \ev{expose}:\ty{Lit}\le\ty{Exact},\qquad
 \ev{forget}:\ty{Exact}\le\ty{Abs}.
\]
The \emph{source} and \emph{destination} of a coercion are its input and
output interfaces, both within the target language. For each destination
proposition, we construct a recipe deriving it from a source proposition.

\subsection*{From a definition to exact bounds}
The source telescope of $\ty{Lit}$ contains one slot:
$q_0:s.A\equiv W$. The destination $\ty{Exact}$ requires two facts:
\[
 \begin{array}{lll}
 \text{required lower bound}&W\le s.A&
       \text{read }q_0\text{ backwards},\\
 \text{required upper bound}&s.A\le W&
       \text{read }q_0\text{ forwards}.
 \end{array}
\]
Equality supports both inclusion directions; inclusion evidence alone
does not support reversal. Both output recipes may therefore use the
same input slot.

The name $q_0$ is a \emph{placeholder} for a source proposition, not a
fact already selected from a particular object. These recipes define a
conversion independently of the object to which it will be applied.
Their ordered list is called a \emph{morphism}.

\subsection*{From exact bounds to abstract bounds}
Now the source slots are
$p_0:W\le s.A$ and $p_1:s.A\le W$.
For the destination lower bound, prepend the inclusion from $\bot$:
\[
 \bot\xrightarrow{\ev{bot}(W)}W
      \xrightarrow{p_0}s.A.
\]
For the destination upper bound, append the inclusion into $\top$:
\[
 s.A\xrightarrow{p_1}W
      \xrightarrow{\ev{top}(W)}\top.
\]
These two recipes discharge the destination obligations using the
source's lower and upper slots, respectively.

\subsection*{Templates and morphisms}
Each of these recipes is represented by a \emph{template} with the three-part form
\[
 \boxed{\text{left side}\quad|\quad\text{source slot}\quad|\quad\text{right side}}.
\]
A side may be absent when its endpoint requires no further conversion.
The central slot is a \emph{hole}, instantiated by a fact about the
object when the coercion is applied. Figure~\ref{fig:recipes} gives all
four recipes and the source subtyping rule translated by $\ev{forget}$.

\begin{figure}[!htbp]
\begin{examplepanel}{srccol}{Source widening --- the declaration subtyping rule}
\[
\src{\inferrule*[right=\textsc{Typ}]{
 \bot<:W\\ W<:\top
}{\iface{W}{W}<:\iface{\bot}{\top}}}
\]
Widening decreases the lower bound and increases the upper bound.
\end{examplepanel}
\begin{examplepanel}{tgtcol}{Target interfaces}
\[
\tgt{\begin{aligned}
 \ty{Lit}&=\obj{s.A\equiv W},\\
 \ty{Exact}&=\obj{W\le s.A,\ s.A\le W},\\
 \ty{Abs}&=\obj{\bot\le s.A,\ s.A\le\top}.
\end{aligned}}
\]
\end{examplepanel}
\begin{examplepanel}{tgtcol}{Coercion templates --- one per destination proposition}
\small
\renewcommand{\arraystretch}{1.3}
\begin{tabularx}{\linewidth}{@{}l l >{\raggedright\arraybackslash}X l@{}}
\toprule
Conversion and output & Left side & Input hole & Right side\\
\midrule
$\ev{expose}$, lower & absent & equality 0, reversed & absent\\
$\ev{expose}$, upper & absent & equality 0, forward & absent\\
$\ev{forget}$, lower & $\ev{bot}(W)$ & inclusion 0 & absent\\
$\ev{forget}$, upper & absent & inclusion 1 & $\ev{top}(W)$\\
\bottomrule
\end{tabularx}
\end{examplepanel}
\caption{The coercion \(\ev{expose}:\ty{Lit}\le\ty{Exact}\) derives
exact bounds from the definition equality, while
\(\ev{forget}:\ty{Exact}\le\ty{Abs}\) translates the displayed
widening derivation. Each hole indexes the source telescope of its
coercion. The templates are independent of the object to which the
coercion will be applied.}
\label{fig:recipes}
\end{figure}


The indices always refer to the \emph{source} telescope of that conversion.
In particular, the two occurrences of index 0 in the first two rows
refer to the same equality.

In compact notation, omitting the constructors that attach a recipe to
its surrounding list, the upper recipe for $\ev{expose}$ is
\[
 \ev{le}\ \ev{none}\ (\ev{eq}\ 0)\ \ev{none}.
\]
The constructor $\ev{le}$ specifies an inclusion output, the two
$\ev{none}$ arguments specify absent sides, and $\ev{eq}\ 0$ reads
equality slot 0 forwards. Replacing $\ev{eq}$ by $\ev{eqSym}$ gives the
lower recipe.

The complete object coercion has the form
$\ev{obj}\ \ty{sourceTelescope}\ m$, where $m$ is the list of recipes.
It carries the source telescope. Its destination is determined by the
facts the recipes establish. The implementation represents these lists
with constructors; the bracket notation here displays their order.

\subsection*{Scope of the template sides}
In a basic template, present sides are scoped outside the self binder
and cannot mention it. For example, $\ev{bot}(W)$ refers only to $W$.
The recipe accesses $s$ through its hole.

The restriction concerns the scope of $s$; an outer variable may still
denote the same runtime object. An absent side can preserve an endpoint
such as $s.A$ without supplying evidence outside the binder.
As Section~\ref{sec:normalization} shows, this representation supports
composition by substitution into holes, without additional member
lookups through self.

\exercise{widen}{Build the conversion from $\ty{Exact}$ to
$\obj{\bot\le s.A,\ s.A\le W}$. Give the lower and upper recipes
in three-part form, and say which side of $\ev{forget}$ you removed.}

\lesson{Stores and views}\label{sec:store}
\focus{A typed store supplies the concrete definitions used to interpret interface evidence.}

The preceding interfaces and coercions can be checked before execution.
To interpret them at runtime, suppose the machine has allocated the
literal at $\ell$:
\[
 \sigma(\ell)=\{A=W\}.
\]
The store records the value itself. Write $\Gamma_\sigma$ for its typing
context: it records the precise type $\ty{Lit}$ and the definition of
$\ell.A$. A binder carrying definitions is called \emph{transparent}.
A function parameter
$y:\ty{Exact}$ is \emph{opaque}: its interface supplies assumptions, but
there is no stored definition for its $y.A$ while checking the body.

The transparent location admits definition evidence
$\ev{def}(\ell,A):\ell.A\equiv W$.
An opaque parameter supplies no definition to justify this constructor.
The coercion $\ev{expose}$ instead uses the equality in its source
telescope, so its construction is independent of the stored definition.

\subsection*{Precise allocation}
Suppose a let binds a literal already cast to $\ty{Exact}$:
\[
 \ev{let}\ y=\cast{\{A=W\}}{\ev{expose}}\ \ev{in}\ t.
\]
The machine stores the bare literal at a fresh location $\ell$ and
substitutes the atom $\cast{\ell}{\ev{expose}}$ for $y$ in the body.
The store therefore retains $\ty{Lit}$ and the definition $A=W$;
the client sees $\ty{Exact}$ through its evidence.

This allocation rule is part of the store invariant. Recording arbitrary
interfaces as store facts would admit assertions unsupported by the
stored values. Instead, the store retains precise definitions, and cast
evidence justifies each change of presentation.

\subsection{Closing the complete source translation}\label{sec:closed-program}
We can now connect the client calculation to the entire running program.
Use $x,f$ for its source binders and reserve $\ell,z$ for the two
locations allocated during execution. This keeps a parameter's assumed
interface distinct from a location's precise stored type.

\textbf{First supply the source object's interface.}
Let $D=\iface{W}{W}$. The source definition rule directly assigns
$A=W$ the exact declaration $D$. The object rule checks that definition
under a self binding at $\mu(s.D)$ that also records the literal's
definitions, and gives the literal type $\mu(s.D)$.
Here $W$ is independent of self, so checking the definition uses no
facts from that binding. After the first let, variable typing gives
$x:\mu(s.D)$, whereas reading a bound needs $x:D[x/s]$.
The source recursive-elimination rule supplies precisely that change:
\[
 \src{\frac{\Gamma\vdash x:\mu(s.D)}
                 {\Gamma\vdash x:D[x/s]}\ \textsc{Rec-E}.}
\]
It changes the derivation of the variable occurrence, not the source
expression. In this example $D$ contains no free occurrence of its
bound self, so $D[x/s]=D$. Nevertheless the rule is present in the
derivation and has a visible translation.

\textbf{Translate introduction and opening.}
The target literal has precise type $\ty{Lit}$, while the translation
of the source literal's assigned type is $\ty{Exact}$. The object
translation therefore inserts the coercion already constructed:
\[
 \tr{d_{\mathrm{obj}}}
       =\cast{\{A=W\}}{\ev{expose}}.
\]
For recursive opening, write
$P_E=[W\le s.A,\ s.A\le W]$ for the telescope of $\ty{Exact}$ and define
the exact wrapper abbreviation
\[
 R(c)=\ev{foldSelf}_{P_E}(\ev{unfoldSelf}(c)).
\]
The unfold wrapper opens the recursive presentation at the atom's root;
the fold wrapper packages the opened declaration's telescope as an
object interface. These are atom constructors, not evidence composition.
Here opening leaves the telescope's external types unchanged, so
$R(c):\ty{Exact}$ whenever $c:\ty{Exact}$.
Both constructors preserve the root. A view records evaluated interface
facts, as the next calculation will explain; their view rules give
$\view{R(c)}=\view{c}$ when interpreted over a typed store.
They are still present in the syntax and consume normalizer fuel.

Thus the bound derivations for this literal-origin source program
translate to $\ev{lower}_A(R(x))$ and $\ev{upper}_A(R(x))$.
The earlier parameter calculation used $\ev{lower}_A(x)$ and
$\ev{upper}_A(x)$ because it assumed $x:D$ directly.
This is the exact difference between the two derivations.

\textbf{Assemble the whole target.}
The identity lambda translates to itself: its body uses variable typing.
Each of the four lets uses the let clause from
Section~\ref{sec:evidence}. With the chosen lower-bound, upper-bound,
and argument subsumptions, the complete translation is
\[
\tgt{\begin{aligned}
 \ev{let}\ x&=\cast{\{A=W\}}{\ev{expose}}\ \ev{in}\\
 \ev{let}\ f&=\lambda(u:\top).u\ \ev{in}\\
 \ev{let}\ \ty{hidden}&=\cast{f}{\ev{lower}_A(R(x))}\ \ev{in}\\
 \ev{let}\ \ty{recovered}
   &=\cast{\ty{hidden}}{\ev{upper}_A(R(x))}\ \ev{in}\\
 &\ty{recovered}\ \cast{f}{\ev{top}(W)}.
\end{aligned}}
\]
Every constructor introduced by this derivation is now accounted for:
the object rule inserts $\ev{expose}$, \textsc{Rec-E} inserts $R$,
the three subsumptions insert casts, and lambda and let translation
retain the program structure. The notation still suppresses the
implementation's variable indices; $R$ and the member-read abbreviations
have explicit definitions.

\textbf{Close the body by allocation and substitution.}
After the first two bindings have allocated their bare values, the store
contains
\[
 \sigma(\ell)=\{A=W\},\qquad
 \sigma(z)=\lambda(u:\top).u,
 \qquad \Gamma_\sigma(\ell)=\ty{Lit},\quad
        \Gamma_\sigma(z)=W.
\]
The literal's wrapper is retained in the atom substituted for $x$:
\[
 a=\cast{\ell}{\ev{expose}},\qquad a^\circ=R(a),\qquad
 \theta=[x:=a,\ f:=z].
\]
For the remaining body $B$ of the displayed target, substitution gives
\[
\tgt{\begin{aligned}
 B\theta={}&
 \ev{let}\ \ty{hidden}=\cast{z}{\ev{lower}_A(a^\circ)}\ \ev{in}\\
 &\ev{let}\ \ty{recovered}
     =\cast{\ty{hidden}}{\ev{upper}_A(a^\circ)}\ \ev{in}\\
 &\ty{recovered}\ \cast{z}{\ev{top}(W)}.
\end{aligned}}
\]
In evidence receivers, $R(x)$ has become the whole atom $R(a)$.
In type selections, $x.A$ has become $\ell.A$, since
$\ev{root}(a)=\ell$. Substituting the two aliases next gives the
closed application
\[
 \cast{\cast{z}{\ev{lower}_A(a^\circ)}}{\ev{upper}_A(a^\circ)}
 \quad \cast{z}{\ev{top}(W)}.
\]
This is the application evaluated in Section~\ref{sec:application}.
The following view calculation supplies its remaining obligation:
the two member reads really normalize to identity at this store.
The companion's \texttt{FullSource} examples check the full source
derivation and its exact target translation; the machine calculation here
uses the stated allocation and substitution rules.


\subsection*{Concrete views}
A \emph{view} records the evaluated information at an atom's object
interface, in telescope order with one entry per proposition.
An inclusion entry contains a \emph{head form}, which records the outer
structure of the inclusion after following stored definitions.
We mark evaluated results with hats. The form $\form{id}$ says the
endpoints agree; $\form{bot}$ says the source endpoint is $\bot$;
$\form{top}$ says the destination endpoint is $\top$.
Thus $\ev{top}(W)$ is evidence to evaluate, whereas $\form{top}$ is its
evaluated form. Section~\ref{sec:normalization} adds function and object
forms and explains their computation.

An equality entry instead contains an equality marker. For our literal:
\[
 \view{\ell}=[\,\form{eq}\,].
\]
The marker $\form{eq}$ says that the equality's endpoints agree after
following the definitions. It does not record a new equality axiom.
The typing of the view relates that marker to the proposition
$\ell.A\equiv W$.

Let
\[
 a=\cast{\ell}{\ev{expose}}:\ty{Exact},\qquad
 b=\cast{a}{\ev{forget}}:\ty{Abs}.
\]

The full translated program uses the reopened receiver
$a^\circ=R(a)$ from Section~\ref{sec:closed-program}. Its wrappers
preserve the view as well as the root:
$\view{a^\circ}=\view{a}$ and $\ev{root}(a^\circ)=\ell$.
We first compute the view of $a$, then use this equality when reading the
program's bounds. Reading those bounds directly from the bare location
$\ell:\ty{Lit}$ would select from the wrong telescope.

To compute $\view{a}$, apply the two $\ev{expose}$ recipes to
$[\,\form{eq}\,]$. Each recipe reads the same equality slot. Forward or
backward, its endpoints resolve to the same type $W$, so each produces
the identity inclusion form:
\[
 \view{a}=[\,\ev{le}\,\form{id},\ \ev{le}\,\form{id}\,].
\]
Here $\ev{le}$ tags an inclusion entry, distinguishing it from an
equality entry.

Applying $\ev{forget}$ uses composition of head forms, written $F;G$.
This operation computes with normalized data, whereas $p;q$ constructs
an evidence expression. The first recipe prepends $\form{bot}$ to the
slot's $\form{id}$; the second appends $\form{top}$:
\[
 \begin{aligned}
 \text{lower output}&:\ \form{bot};\form{id}=\form{bot},\\
 \text{upper output}&:\ \form{id};\form{top}=\form{top}.
 \end{aligned}
\]
Thus
\[
 \view{b}=[\,\ev{le}\,\form{bot},\ \ev{le}\,\form{top}\,].
\]
The same literal now supplies the weaker facts of $\ty{Abs}$.
From that interface alone, the client cannot recover $W$;
the store still records that definition. Figure~\ref{fig:presentations}
records the three views and the normalization of the running program's
member evidence. There $\hnf{e}$ denotes head normalization of evidence,
the computation that produces its head form.

\needspace{4\baselineskip}
The notation $\operatorname{resolve}_{\Gamma_\sigma}(T)$ means following
the definitions in $\Gamma_\sigma$ to expose the outer shape of $T$.
Here it turns $\ell.A$ into $W$.


\begin{figure}[!htbp]
\begin{examplepanel}{runcol}{Stored definition and type resolution}
\[
 \run{\sigma(\ell)=\{A=W\},\qquad
       \operatorname{resolve}_{\Gamma_\sigma}(\ell.A)=W.}
\]
\end{examplepanel}
\begin{examplepanel}{tgtcol}{Three target atoms and their concrete views}
\small
\renewcommand{\arraystretch}{1.35}
\begin{tabularx}{\linewidth}{@{}l >{\raggedright\arraybackslash}X l@{}}
\toprule
Interface & Atom rooted at \(\ell\) & View\\
\midrule
$\ty{Lit}$ & $\ell$ & $[\,\form{eq}\,]$\\
$\ty{Exact}$ & $a=\cast{\ell}{\ev{expose}}$
 & $[\,\ev{le}\,\form{id},\,\ev{le}\,\form{id}\,]$\\
$\ty{Abs}$ & $b=\cast{a}{\ev{forget}}$
 & $[\,\ev{le}\,\form{bot},\,\ev{le}\,\form{top}\,]$\\
\bottomrule
\end{tabularx}
\end{examplepanel}
\begin{examplepanel}{tgtcol}{Normalization of the running program's member evidence}
\[
\tgt{\begin{aligned}
 \ev{lower}_A(a^\circ)&:W\le \ell.A
     &\hnf{\ev{lower}_A(a^\circ)}&=\form{id},\\
 \ev{upper}_A(a^\circ)&:\ell.A\le W
     &\hnf{\ev{upper}_A(a^\circ)}&=\form{id}.
\end{aligned}}
\]
\end{examplepanel}
\caption{Applying coercion templates to the stored object's facts
produces concrete views. The atom \(a:\ty{Exact}\) instantiates the
allocated client binding, and $a^\circ=R(a)$ is its reopened
member receiver. Their root \(\ell:\ty{Lit}\) supplies the dependent
name \(\ell.A\). The abstract presentation exposes weaker interface facts
while the store retains the definition.}
\label{fig:presentations}
\end{figure}



\remember{Evidence expressions, normalized object recipes, and concrete
views are distinct representations. Normalized recipes retain holes;
applying them to a stored object's facts produces a concrete view.}

\exercise{views}{Through $a$, what head form does reading the upper bound
produce? What about through $b$? Explain why these answers can differ
even though $\ev{root}(a)=\ev{root}(b)=\ell$.}

\lesson{Evidence normalization and composition}\label{sec:normalization}
\focus{Head forms support composition while retaining the evidence needed for subsequent operations.}

The three head forms used in the view calculation record identity,
an inclusion from $\bot$, or an inclusion into $\top$.
Function and object conversions need forms that retain their component
evidence or templates. For our fragment, the possibilities are:
\begin{center}
\renewcommand{\arraystretch}{1.2}
\begin{tabularx}{\linewidth}{@{}lX@{}}
\toprule
Form & What its typing establishes after resolving definitions\\
\midrule
$\form{bot}$ & The source endpoint is $\bot$.\\
$\form{top}$ & The target endpoint is $\top$.\\
$\form{id}$ & The two endpoints agree.\\
$\form{pi}(d,c)$ & Both endpoints are functions; $d,c$ adapt argument and result.\\
$\form{obj}(E_0,\ldots,E_n)$ & The object conversion has these normalized recipes.\\
\bottomrule
\end{tabularx}
\end{center}
This table covers the forms used in the running examples. The full
target also has forms for self-bounds, discussed in
Section~\ref{sec:larger}.

The form $\form{id}$ carries no endpoints; its typing judgment requires
them to resolve equally. It is therefore a valid form for $\ell.A\le W$
when the store defines $\ell.A$ as $W$.

\subsection*{Normalizing the round trip}
The view at $\ty{Exact}$ gives
\[
 \begin{aligned}
 \hnf{\ev{lower}_A(a^\circ)}&=\form{id},\\
 \hnf{\ev{upper}_A(a^\circ)}&=\form{id}.
 \end{aligned}
\]
For composition the normalizer combines the two head forms,
using the operation on forms introduced in the view calculation:
\[
 \hnf{\ev{lower}_A(a^\circ);\ev{upper}_A(a^\circ)}
   =\form{id};\form{id}=\form{id}.
\]
The original inclusions have endpoints $W\le \ell.A$ and $\ell.A\le W$,
and their composite has endpoints $W\le W$. Resolution justifies the
identity forms at the intermediate endpoints.

\subsection*{Composing interface coercions}
The coercions $\ev{expose}$ and $\ev{forget}$ can also be composed
before either is applied to $\ell$.
The lower recipe of $\ev{forget}$ is
\[
 [\,\ev{bot}(W)\mid\text{inclusion slot }0\mid\text{absent}\,].
\]
Slot 0 of its input is the lower recipe produced by
$\ev{expose}$:
\[
 [\,\text{absent}\mid\text{equality slot }0\text{ reversed}\mid
       \text{absent}\,].
\]
Substituting the second recipe into the first hole yields
\[
 [\,\ev{bot}(W)\mid\text{equality slot }0\text{ reversed}\mid
       \text{absent}\,].
\]
The upper composite is
\[
 [\,\text{absent}\mid\text{equality slot }0\text{ forward}\mid
       \ev{top}(W)\,].
\]
The resulting morphism converts $\ty{Lit}$ directly to $\ty{Abs}$.
Normalizing its sides leaves the equality holes intact; instantiating
them with $\view{\ell}$ gives $\form{bot}$ and $\form{top}$, as with the
two successive view transformations. The morphism representation thus
supports composition by an explicit substitution operation on finite
templates.

\subsection*{Scope of head normalization}
The representation determines how far normalization proceeds.
Basic object templates access self through designated slots, with side
evidence outside the self binder. Composition can therefore combine
existing components without introducing a new member lookup through self.

A function form $\form{pi}(d,c)$ retains its argument and result
adaptations as evidence. In the dependent case, $c$ may mention a parameter
whose value has not yet been supplied. The machine needs only the outer
function form to perform a call; once the argument is available, it uses
$d$ and $c$ as casts. Head normalization therefore does not require
store-based normalization under an unapplied parameter.

The executable normalizer in Lean takes \emph{fuel}: a natural-number
budget that decreases during recursion. It terminates on every input
and returns an optional form, with failure possible both for malformed
evidence and for insufficient fuel. Failure at a particular budget
therefore does not establish that the evidence is ill typed.

The normalization theorem applies to evidence \emph{closed over a typed
store}: every free variable names a location in that store, whose value
has been checked at its precise type. Such evidence may mention $\ell$;
its checking context contains no extra opaque parameter assumptions.
For well-typed evidence in this setting, there is a fuel budget for which
normalization succeeds with a form of the required type. The proof is by
structural induction on typing derivations.
This is the inclusion-evidence part of the \emph{canonical-forms theorem}.

\exercise{compose}{Fill the slot in
$[\,\ev{bot}(W)\mid\text{equality 0 reversed}\mid\text{absent}\,]$
using the concrete view $[\,\form{eq}\,]$. Which head form results?
Why is the unfilled recipe not already a concrete view?}

\lesson{Function coercions and application}\label{sec:application}
\focus{A function head form supplies the argument and result adaptations required by an application.}

In the store constructed in Section~\ref{sec:closed-program}, location
$z$ holds $\lambda(u:\top).u$ at precise type $W=\top\to\top$.
Using the reopened exact-interface atom $a^\circ=R(a)$, that section's
closing calculation gives
\[
 z'=\cast{\cast{z}{\ev{lower}_A(a^\circ)}}{\ev{upper}_A(a^\circ)}:W.
\]
The chain of casts from $z$ to $z'$ normalizes to $\form{id}$.
If $r:\top$ is an argument atom, the machine can therefore call the
lambda stored at $z$ directly. Substitution in the body $u$ returns $r$.
For the original program, choose
$r=\cast{z}{\ev{top}(W)}$.
The result erases to $z$, as expected from evaluation of the identity.

Establishing this step requires more than the type $z':W$: the atom
$z'$ is a cast expression, while application must locate a stored lambda.
The canonical-forms theorem connects the atom's presented type to the
value at its root.

\subsection*{One machine step, with its premises}
Write $\langle\sigma,K,t\rangle$ for a state with store $\sigma$,
continuation $K$, and current term $t$. Let
$\operatorname{chain}_\sigma(c)$ normalize the wrappers from an atom's
root to its presented type. For a wrapped receiver, the identity case
of the machine's application rule is
\[
 \frac{\sigma(\ev{root}(c))=\lambda(u:S).t_0
       \qquad c\ne\ev{root}(c)
       \qquad \operatorname{chain}_\sigma(c)=\form{id}}
      {\langle\sigma,K,c\,r\rangle
         \longrightarrow\langle\sigma,K,t_0[u:=r]\rangle}.
\]
The comparison in the middle treats the root as a bare atom.
The chain equation abbreviates successful normalization at some fuel;
it is a premise of the step, not an unchecked optimization.

For the closed program, set $c=z'$, $K=\epsilon$, $S=\top$, and
$t_0=u$. The root is $z$, its stored value is the identity lambda,
and both member casts have identity form by the preceding view
calculation. Hence all three premises hold, and the actual step is
\[
 \langle\sigma,\epsilon,z'\,\cast{z}{\ev{top}(W)}\rangle
 \longrightarrow
 \langle\sigma,\epsilon,\cast{z}{\ev{top}(W)}\rangle.
\]
This is \texttt{Step.appCastRefl}. It changes neither the store nor the
continuation, and its result is an atom. Earlier allocation and alias
substitution took several steps; this display isolates one complete
transition and all the evidence needed to enable it.

\subsection*{Adapting an argument}
For a nonidentity function coercion, present $z$ at the function type
$W\to\top$, whose domain is narrower than that of $W$.
The required inclusion is
\[
  (\top\to\top)\le(W\to\top).
\]
Clients of the new interface supply arguments at $W$, which the original
function can accept at $\top$. The argument and result evidence are
\[
 d=\ev{top}(W):W\le\top,\qquad
 c=\ev{refl}(\top):\top\le\top.
\]
The function constructor builds
$\ev{pi}(d,c):(\top\to\top)\le(W\to\top)$.
Head normalization exposes its function form:
\[
 \hnf{\ev{pi}(d,c)}=\form{pi}(d,c).
\]
The components $d,c$ remain evidence for the call's argument and result
casts. The domain adaptation is contravariant: the function inclusion
maps the original function interface to the client interface, while
$d$ maps the client's argument type to the original domain.

Figure~\ref{fig:function-call} gives the source rule, its target
evidence, and the call. A client supplies $z:W$; the argument evidence
presents it at $\top$ for the stored identity. The body returns that
same atom, and the result evidence presents it to the client.
Administrative steps arrange the pending casts around evaluation of the
function body. The call erases to $z\,z$ and its result to $z$.


\begin{figure}[!htbp]
\begin{examplepanel}{srccol}{Source function subtyping --- contravariant domain}
\[
\src{\inferrule*[right=\textsc{All}]{
 W<:\top\\ \top<:\top
}{(\top\to\top)<:(W\to\top)}}
\]
\end{examplepanel}
\begin{examplepanel}{tgtcol}{Target evidence and application}
\[
\tgt{\begin{aligned}
 d&=\ev{top}(W):W\le\top,&c&=\ev{refl}(\top):\top\le\top,\\
 g&=\cast{z}{\ev{pi}(d,c)}:W\to\top.
\end{aligned}}
\]
\par\small
\renewcommand{\arraystretch}{1.2}
\noindent\begin{tabularx}{\linewidth}{@{}lX@{}}
\toprule
Stage of \(g\,z\) & Argument or result at that stage\\
\midrule
Client supplies & $z:W$\\
Original function receives & $\cast{z}{d}:\top$\\
Identity body returns & $\cast{z}{d}:\top$\\
Client receives & $\cast{\cast{z}{d}}{c}:\top$\\
\bottomrule
\end{tabularx}\par
\end{examplepanel}
\begin{examplepanel}{runcol}{Erased execution}
\[
\run{\erase{\sigma}(z)=\lambda u.u,\qquad z\,z\longrightarrow z.}
\]
\end{examplepanel}
\caption{The function coercion adapts the argument with \(d\) and the
result with \(c\). The typed store holds \(z=\lambda(u:\top).u\), and
the function head form retains both evidence components for application.
The table groups administrative steps; the erased call applies the
stored identity.}
\label{fig:function-call}
\end{figure}



\FloatBarrier
\needspace{8\baselineskip}
\subsection*{Dependent results and composition}
For nondependent function types, $c$ is independent of the argument.
For a dependent function $(u:S)\to T(u)$, the result type can mention
$u.A$. At a call, substitute the argument atom in term positions and its
root in type positions. A wrapper changing the argument's interface
must not create a different dependent result name.

Composition of function coercions respects this order of adaptation.
If a first adaptation expects an intermediate input and a second offers
a still narrower one, the second input adaptation runs first.
Result adaptations run in the order of the function coercions.
In a dependent codomain, the earlier result evidence is instantiated
with the adapted argument before being composed with the later result evidence.

\remember{Application uses a function head form to adapt the argument,
evaluate the stored lambda, and adapt the result. The component evidence
is retained until the argument is available.}

\exercise{variance}{Suppose $d:S_2\le S_1$ and $c:T_1\le T_2$,
with nondependent results. Give the function inclusion built from them.
When it is applied to $b:S_2$, at which type does the original body
receive $b$, and at which type does the client receive the result?}

\lesson{Bad bounds and uninhabited interfaces}\label{sec:bad}
\focus{Inconsistent bounds can occur in an open context, but no typed store supplies an argument satisfying them.}

Consider the source interface $\iface{\top}{\bot}$ and its target
translation
\[
 \ty{Bad}=\obj{\top\le s.A,\ s.A\le\bot}.
\]
Under a hypothetical parameter $y:\ty{Bad}$, we can read both facts
and compose them:
\[
 \ev{lower}_A(y);\ev{upper}_A(y):\top\le\bot.
\]
This evidence is well typed in the open context. In particular, a
function body may be checked under an uninhabited parameter type.
The safety argument must show that execution cannot supply an argument
at that type.

\needspace{10\baselineskip}
\subsection*{The obstruction to constructing an argument}
Our stored object defines $A=W$.
To present it at $\ty{Bad}$ using the recipe pattern, we would need
both chains:
\[
 \begin{aligned}
 \text{lower: }&\top\le W\le \ell.A,\\
 \text{upper: }&\ell.A\le W\le\bot.
 \end{aligned}
\]
The definition equality handles the links next to $\ell.A$.
The remaining obligations are $\top\le W$ and $W\le\bot$.
Composing a successful construction would produce closed evidence for
$\top\le\bot$ over the typed store.

The normalization theorem excludes such evidence. None of the typed
head forms for this fragment admits those endpoints:
$\top$ is not $\bot$; $\bot$ is not $\top$; they do not resolve equally;
they are not both functions; and they are not both object interfaces.
The full target's consistency result covers its additional forms as well.

The same consistency result excludes every atom at $\ty{Bad}$ over a
typed store: selecting and composing its bounds would again yield
closed evidence for $\top\le\bot$.

\subsection*{Open assumptions and store contexts}
The open derivation used $y:\ty{Bad}$ as an assumption.
The runtime theorem requires a store context, which records the stored
values' precise types and definitions. Well-scopedness or well-formedness of
an open context alone does not establish this hypothesis.

At application, substitution connects the open body, checked under its
opaque parameter, to a typed atom over the store. No such atom exists
for this parameter type, so execution cannot reach the body through a
well-typed call. The argument does not require normalizing the open
evidence under the impossible assumption.

\exercise{bad}{Consider the claim that the well-typed open evidence
$\ev{lower}_A(y);\ev{upper}_A(y)$ must have a typed head form, proving
the system inconsistent. Identify the missing theorem hypothesis precisely.}

\lesson{Canonical forms and type safety}\label{sec:proof}
\focus{The canonical-forms theorem connects evidence normalization to the operations of the target machine.}

The proof establishes properties of several judgments simultaneously,
following the dependencies in the syntax: cast atoms contain inclusion evidence, member selection
requires an atom, and object coercions contain templates. Normalization
of inclusion evidence therefore depends on properties of atoms and views.

\needspace{15\baselineskip}
\subsection*{The canonical-forms theorem}
Fix a typed store $\sigma$ with its context $\Gamma_\sigma$.
All judgments below are in that store context, so their free variables
are allocated locations. The canonical-forms theorem establishes the
following related facts:
\begin{enumerate}
\item Well-typed inclusion evidence has a typed head form.
\item Equality evidence relates endpoints that resolve equally.
\item A well-typed atom has a typed view at its object interface.
\item The atom's chain of wrappers has a form relating the root's
precise type to its presented type.
\item For the full target, field-presence evidence identifies a field actually stored in
the appropriate object.
\end{enumerate}
The proof also excludes atoms whose types resolve to $\bot$. Alongside
successful normalization, these statements establish the typing and
semantic properties of the resulting forms.

Three representative cases illustrate the dependencies. For transitivity,
the induction hypotheses
provide typed forms for both components; a composition lemma shows that
combining them preserves typing. For a member read, the hypotheses provide the atom's
view and the interface conversion's normalized recipes; selecting the
appropriate resulting entry gives the needed fact. For a cast atom,
the proof transforms its existing view and composes its chain.
Each case uses an operation from the preceding calculations.

\needspace{8\baselineskip}
\subsection*{Target progress and preservation}
A machine state consists of a store, a current expression, and a
continuation recording pending lets and casts.

\emph{Progress} says a well-typed machine state is final or can step.
At application, the canonical-forms facts identify a lambda at the
function atom's root and give the required argument and result evidence.
At field projection, presence evidence identifies a field to read.
The remaining cases handle lets and casts.

\emph{Preservation} says a step keeps the program and store typed,
allowing for newly allocated locations. Substitution must transport
evidence as well as terms. Allocation must keep the bare value at its
precise type and move its presentations into the body. A function call
must cast the argument and result at the endpoints promised by its form.

Together, progress and preservation ensure that a closed well-typed
target program never reaches a stuck state. The program may diverge;
termination of evidence normalization does not imply termination of
program execution.

\subsection*{Transporting safety to the source}
The translation takes a source \emph{typing derivation}.
For a well-formed source context, it produces a well-typed target term.
Different source derivations can insert different evidence,
but the translation erases to the same underlying source program:
\[
 \erase{\operatorname{translate}(h)}=\erase{t}
 \quad\text{when }h\text{ derives }\Gamma\vdash t:T.
\]
Transporting safety additionally requires operational correspondence.

Both source and target erase into a shared runtime. A simulated source
state is one for which a typed target state has the same erasure.
Translation establishes this relation initially. When the source steps,
runtime correspondence and target simulation provide a matching target
execution; preservation re-establishes the relation.
Conversely, target progress, together with reflection of runtime steps
by the source, shows the source state is final or can step.
Here \emph{reflection} means that a runtime step out of an erased source
state comes from a corresponding source step. Pending administrative
cast steps can be cleared in finitely many steps before making this
argument. Thus no infinite execution can consist solely of administrative
steps. These additional target steps explain
why the two machines do not run in lockstep.

\remember{Canonical forms establish the runtime structure needed for
target safety. Translation and operational correspondence then transport
safety to the source: erasure relates program syntax, and simulation
relates executions.}

For the main source calculus, this argument derives safety from the
target's metatheory and the simulations, without a separate source
preservation proof as a premise.

\exercise{proof}{Which fact is missing from each proposed proof?
(a) A well-typed translation is claimed to establish source safety by itself.
(b) The existence of one target step is claimed to establish safety for
all future steps. Give a different missing ingredient for each.}

\lesson{Explicit evidence for subcapturing}\label{sec:capture}
\focus{Subcapturing evidence relates capture sets through their interpretation as capability roots.}

The core argument checked evidence under assumptions, then closed it over
a typed store before interpreting it. The capture extension adds another
kind of evidence to the same process. Its interpretation will bound
capability roots; substituting a call's argument into that evidence will
connect a closure's advertised captures to the uses of its executing body.

A capturing type $S^C$ pairs a value's shape $S$ with a capture set $C$
describing its associated capabilities. In the examples,
$\ty{Unit}=\top^\varnothing$ abbreviates the pure universal shape, and
$\ev{unit}$ is a stored pure value of that type. This abbreviation uses
the existing type language without adding a primitive unit type.

Assume the platform supplies two distinct rigid capability roots,
$\kappa_{\ell}$ and $\kappa_c$. A rigid root denotes itself; a program
name is interpreted through its capture annotation. The store contains
closures $\ev{log}$ and $\ev{console}$ with annotations
$\{\kappa_{\ell}\}$ and $\{\kappa_c\}$, respectively. Both are identity
closures on $\ty{Unit}$, with nonempty annotations that serve as
conservative allowances. Their names indicate the intended capability
roles; logging and console I/O are outside the operational semantics.

\subsection*{Subcapturing and root interpretation}
The sets $\{\ev{log}\}$ and $\{\kappa_{\ell}\}$ contain different names,
so they are unrelated by ordinary subset inclusion. Subcapturing also
uses the annotation of $\ev{log}$, which accounts for its capabilities
through $\kappa_{\ell}$:
\[
 \ev{capvar}(\ev{log}):
       \{\ev{log}\}\le_c\{\kappa_{\ell}\}.
\]
In general, if a program name $x$ has type $S^C$, then
$\ev{capvar}(x)$ proves $\{x\}\le_c C$.
The constructor $\ev{elem}$ represents ordinary subset inclusion:
if every element of $C$ is listed in $D$, then
$\ev{elem}(C,D):C\le_c D$.
Transitivity composes capture evidence as it does type evidence.
Union combines evidence for two sets with a shared destination:
from $p:C\le_c D$ and $q:C'\le_c D$, the expression
$\ev{union}(p,q)$ proves $C\cup C'\le_c D$.

Capture evidence is interpreted through \emph{roots}. For the acyclic
annotations used here, $\roots{C}$ follows the capture information
associated with the names in $C$ to their underlying capability roots:
\[
 \roots{\{\ev{log}\}}=\{\kappa_{\ell}\},\qquad
 \roots{\{\ev{console}\}}=\{\kappa_c\},\qquad
 \roots{\{\ev{unit}\}}=\varnothing.
\]
Thus $\roots{C}$ is a set of capability roots, while $\ev{root}(a)$ is
the single program variable obtained by stripping an atom's wrappers.

Over a typed store, well-typed capture evidence has the interpretation
\[
  e:C\le_c D\quad\Longrightarrow\quad
  \roots{C}\subseteq\roots{D}.
\]
This canonical-forms result depends on the typed-store hypothesis.
As with type members, capture members can supply assumptions whose
semantic interpretation requires an actual stored value.

\subsection*{Immediate use sets}
The use-set function assigns a syntactic use bound to each term.
The cases needed for these examples are:
\[
\tgt{\begin{aligned}
\uses{a}&=\{\ev{root}(a)\},& \uses{v}&=\varnothing,\\
 \uses{a\,b}&=\{\ev{root}(a),\ev{root}(b)\},&
 \uses{\ev{let}\ x=t\ \ev{in}\ u\langle U;q\rangle}
   &=\uses{t}\cup U.
\end{aligned}}
\]
Here $v$ denotes a value expression, such as a lambda literal, whose
creation leaves its body unevaluated. A stored name $x$ is an atom with
use set $\{x\}$; its roots are determined by its annotation. Thus the
returned name $\ev{unit}$ has a singleton use set and no capability
roots. The term \emph{immediate} excludes the latent bodies of value
expressions while still including pending let bodies.

In the let's annotation $\langle U;q\rangle$, $U$ is an upper bound
for the pending body's uses, expressed using names available outside
the let. In the body's context, where $x$ is available, $q$ proves
$\uses{u}\le_c U$. Replacing a local name by outer names that account
for it is called \emph{avoidance}. The factory example below includes
an explicit construction of this evidence.

An atom's wrappers affect its typing but preserve its root. In particular,
$\uses{\cast{a}{e}}=\{\ev{root}(a)\}$ when the cast is an atom,
and a term cast satisfies $\uses{\cast{t}{e}}=\uses{t}$.
Evidence can mention names without adding them to immediate uses:
this function bounds the program's modeled inspections, not all names
appearing in its syntax. These clauses matter when an argument is a
wrapped atom rather than a bare variable.

\subsection*{Translating closing evidence}
The source judgment $\src{C;\Gamma\vdash t:T}$ assigns type $T$ to $t$
with permitted immediate uses $C$. A source lambda contains its parameter
type and body; its typing derivation selects the capture annotation $A$
and checks the body with allowance $A\cup\{u\}$. Translation makes this
annotation and its closing evidence explicit in the target lambda.

In Figure~\ref{fig:capture-translation}, translation brackets also apply
to capture sets, types, and contexts. The source context $\Gamma$ is
assumed well formed, and $\mathrm{wf}$ denotes well-formedness.
The judgment symbol $\vdash^c$ denotes capture-evidence typing. The superscript
$\uparrow$ preserves references to outer bindings under a new parameter
binder. With distinct printed names, those names keep their spelling;
the named calculations below leave this shift implicit.

The translation computes two components from the body derivation: a
target body and evidence that its computed use set lies within the source
allowance. The target lambda stores both components, just as the
translation of source subsumption records inclusion evidence in a cast.


\begin{figure}[!htbp]
\begin{examplepanel}{srccol}{Source lambda rule}
\small
\[
\src{\inferrule*[right=\textsc{All-I}]{
 d_b:A^\uparrow\cup\{u\};\Gamma,u:T\vdash t:T'\\
 T\ \mathrm{wf}
}{\varnothing;\Gamma\vdash\lambda(u:T).t:((u:T)\to T')^A}}
\]
Only the parameter type and body occur in the source lambda. Its typing
derivation selects the capture annotation and records the body's permitted uses.
\end{examplepanel}
\begin{examplepanel}{tgtcol}{Target lambda with closing evidence}
\[
\tgt{
 \lambda^{\tr{A}}(u:\tr{T}).\
       \underbrace{\tr{d_b}}_{\text{body}}\
       \left\langle
       \underbrace{\tr{d_b}_{\mathrm{uses}}}_{\text{closing evidence }g}
       \right\rangle .}
\]
\small
\[
\tgt{
\tr{\Gamma},u:\tr{T}\vdash^c
 g:\uses{\tr{d_b}}\le_c\tr{A}^{\uparrow}\cup\{u\}.}
\]
\end{examplepanel}
\caption{Translation of a lambda premise in a well-formed source context
\(\Gamma\). The two components are the translated body \(\tr{d_b}\)
and its closing evidence \(\tr{d_b}_{\mathrm{uses}}\).
The superscript \(\uparrow\) carries outer names under the new parameter
binder. Angle brackets delimit stored evidence; the lambda receives only
its term argument at runtime. The functions \texttt{translate} and
\texttt{translateUses} construct these components.}
\label{fig:capture-translation}
\end{figure}


\FloatBarrier
\subsection*{Computing a closing certificate from a derivation}
The translation contract in Figure~\ref{fig:capture-translation} tells us
what $g$ proves. To see how it is built, translate the source body
$\ev{log}\ u$ rather than choosing its certificate by hand.
Write $F=\ty{Unit}\to\ty{Unit}$ and $U=\{\ev{log},u\}$.
All the following judgments are under
$\Delta=\Gamma,u:\ty{Unit}$, where
$\Gamma(\ev{log})=F^{\{\kappa_\ell\}}$.

The source variable rule records the variable itself in both its use set
and its assigned capture set:
\[
 \src{\{\ev{log}\};\Delta\vdash\ev{log}:F^{\{\ev{log}\}}},
 \qquad
 \src{\{u\};\Delta\vdash u:\top^{\{u\}}}.
\]
In particular, the second conclusion is not yet the argument type
$\ty{Unit}=\top^\varnothing$. Source subsumption can change
\emph{both} an assigned type and a permitted use set:
\[
 \src{\frac{d:C;\Delta\vdash t:S\qquad e:\Delta\vdash S<:T
                    \qquad f:\Delta\vdash C\le_c D}
           {\textsc{Sub}(d,e,f):D;\Delta\vdash t:T}}.
\]
These are independent obligations. For the type change on $u$, its
declared pure type justifies $\{u\}\le_c\varnothing$ by the source
variable subcapturing rule. For its use allowance, ordinary inclusion
justifies $\{u\}\le_c U$. For $\ev{log}$, we similarly recover its
declared type using $\{\ev{log}\}\le_c\{\kappa_\ell\}$, and widen its
permitted uses using $\{\ev{log}\}\le_c U$.

Let $d_\ell$ and $d_u$ be these two subsumed variable derivations.
Both have allowance $U$, so the source application rule combines them:
\[
 \src{\frac{d_\ell:U;\Delta\vdash\ev{log}:F^{\{\kappa_\ell\}}
                   \qquad d_u:U;\Delta\vdash u:\ty{Unit}}
           {d_b=\textsc{App}(d_\ell,d_u):
                      U;\Delta\vdash\ev{log}\ u:\ty{Unit}}}.
\]
This specifies the derivation whose translation we will compute.

\textbf{The body component.}
Two target constructors record the capture information just used.
$\ev{recap}(a,p)$ preserves an atom's shape and assigns capture set
$D$ when $p:\{\ev{root}(a)\}\le_c D$.
$\ev{capt}(e,p)$ combines shape inclusion $e$ with capture inclusion
$p$ to convert a capturing type. Thus variable translation and the two
type-subsumption premises give the exact atoms
\[
\tgt{\begin{aligned}
 L&=\cast{\ev{recap}(\ev{log},\ev{refl}\{\ev{log}\})}
          {\ev{capt}(\ev{refl}(F),\ev{capvar}(\ev{log}))},\\
 B&=\cast{\ev{recap}(u,\ev{refl}\{u\})}
          {\ev{capt}(\ev{refl}(\top),\ev{capvar}(u))}.
\end{aligned}}
\]
Here reflexivity on a capture set is capture evidence.
The $\ev{capvar}$ expressions read the declared types of the bare
bindings. They are the translations of the source variable-subcapturing
premises. The resulting body is $\tr{d_b}=L\,B$, with
$\uses{L\,B}=\{\ev{log},u\}$ because both wrappers preserve roots.

\textbf{The use-evidence component.}
Only three clauses of $\tr{-}_{\mathrm{uses}}$ are needed:
\[
\tgt{\begin{aligned}
 \tr{\textsc{Var}(x)}_{\mathrm{uses}}&=\ev{refl}\{x\},\\
 \tr{\textsc{Sub}(d,e,f)}_{\mathrm{uses}}
     &=\tr{d}_{\mathrm{uses}};\tr{f},\\
 \tr{\textsc{App}(d_1,d_2)}_{\mathrm{uses}}
     &=\ev{union}(\tr{d_1}_{\mathrm{uses}},
                  \tr{d_2}_{\mathrm{uses}}).
\end{aligned}}
\]
The type-subtyping premise $e$ affects the body cast; the distinct
use-subcapturing premise $f$ affects this certificate.
Applying the clauses from the variable leaves upwards gives
\[
\tgt{\begin{aligned}
 p_\ell&=\ev{refl}\{\ev{log}\};
                  \ev{elem}(\{\ev{log}\},U),\\
 p_u&=\ev{refl}\{u\};\ev{elem}(\{u\},U),\\
 g_{\mathrm{gen}}&=\ev{union}(p_\ell,p_u):
                     \{\ev{log},u\}\le_c U.
\end{aligned}}
\]
The two branches share destination $U$, exactly as required by
$\ev{union}$. Finally, lambda translation stores both computed components:
\[
 \tgt{\lambda^{\{\ev{log}\}}(u:\ty{Unit}).
                 L\,B\ \langle g_{\mathrm{gen}}\rangle.}
\]
This is the generated syntax for the specified source derivation, in
named notation. The simpler certificate $\ev{elem}(U,U)$ has the same
endpoints but is a different expression. We use that hand-written
certificate in the next trace to keep the operational calculation short.
The substitution argument after the trace applies to either certificate.


\needspace{14\baselineskip}
\subsection*{Closing a closure body}
To isolate the execution argument, consider the simpler hand-written
target closure
\[
 h=\lambda^{\{\ev{log}\}}(u:\ty{Unit}).\,\ev{log}\ u.
\]
Its superscript records the capture annotation $\{\ev{log}\}$, while
its body has immediate use set $\{\ev{log},u\}$. Since the argument is
available when the closure is called, the closing obligation is
\[
 g:\uses{\ev{log}\ u}\le_c\{\ev{log}\}\cup\{u\}.
\]
The two endpoints are identical, so $\ev{elem}$ supplies $g$.
This shares the generated closure's use-set obligation, while its body
and certificate omit the translation's additional syntax.

Figure~\ref{fig:capture-call} gives the closing certificate, root
calculation, and complete two-call trace for this hand-written target
example. A body produced by source translation may contain additional
atom wrappers, with generated closing evidence built from unions and
compositions.


\begin{figure}[!htbp]
\begin{examplepanel}{tgtcol}{Stored closure and closing certificate}
\[
\tgt{\begin{aligned}
 h&=\lambda^{\{\ev{log}\}}(u:\ty{Unit}).\
          \ev{log}\ u\ \langle g_h\rangle,\\
 g_h&=\ev{elem}(\{\ev{log},u\},\{\ev{log},u\}),\\
 g_h&:\uses{\ev{log}\ u}\le_c\{\ev{log}\}\cup\{u\}.
\end{aligned}}
\]
The store also contains \(\ev{unit}:\ty{Unit}\) and the identity closure
\(\ev{log}:(\ty{Unit}\to\ty{Unit})^{\{\kappa_\ell\}}\).
\end{examplepanel}
\begin{examplepanel}{runcol}{Capture-annotation resolution}
\[
\run{\{h\}\ \rightsquigarrow\ \{\mathsf{log}\}\
   \rightsquigarrow\ \{\kappa_\ell\},
   \qquad \operatorname{roots}_{\sigma}(\{\mathsf{unit}\})=\varnothing.}
\]
The relation \(\rightsquigarrow\) resolves capture annotations.
\end{examplepanel}
\begin{examplepanel}{runcol}{Execution with an empty continuation}
\small
\renewcommand{\arraystretch}{1.3}
\begin{tabularx}{\linewidth}{@{}lXll@{}}
\toprule
Current term & Use set & Its roots & Inspection\\
\midrule
$h\ \mathsf{unit}$ & $\{h,\mathsf{unit}\}$ & $\{\kappa_\ell\}$ & $h$\\
$\mathsf{log}\ \mathsf{unit}$ & $\{\mathsf{log},\mathsf{unit}\}$ & $\{\kappa_\ell\}$ & $\mathsf{log}$\\
$\mathsf{unit}$ & $\{\mathsf{unit}\}$ & $\varnothing$ & none\\
\bottomrule
\end{tabularx}
\end{examplepanel}
\caption{Closing evidence, root resolution, and execution of the closure.
The two applications inspect \(h\) and \(\mathsf{log}\), then return an
atom with a singleton use set and no capability roots.}
\label{fig:capture-call}
\end{figure}



Each application inspects the closure named in the last column: the
first executes $h$'s body and the second executes the identity body of
$\ev{log}$, returning the atom $\ev{unit}$. The syntactic use sets change
along this execution while their roots remain bounded by the initial root set.

\FloatBarrier
\needspace{12\baselineskip}
\subsection*{Capture prediction for application}
\label{sec:app-prediction}
We now derive the application bound for an arbitrary argument atom,
including its wrappers. Fix a typed store $\sigma:\Gamma_\sigma$ with
$\sigma(h)=\lambda^A(u:T).t\langle g\rangle$, and suppose
$\Gamma_\sigma\vdash b:T$. The argument must have the stored lambda's
domain type; a coerced call first supplies the required domain cast.
For the present derivation the function receiver is the bare location
$h$, and the continuation is empty.

\textbf{1. Recover the open certificate.}
Typing of the stored lambda gives
\[
 \Gamma_\sigma,u:T\vdash^c
       g:\uses{t}\le_c A^\uparrow\cup\{u\}.
\]
The body and its certificate are still open in $u$. Store typing is what
connects this certificate to the lambda actually found at $h$; an
annotation on a hypothetical parameter alone would not supply that body.

\textbf{2. Substitute the typed argument.}
Put $r=\ev{root}(b)$ and write $\theta=[u\mapsto r]$ for substitution
on capture sets. Term and evidence substitution replace term-level
occurrences of $u$ by the whole atom $b$; their type and capture names
use $r$. The capture-evidence substitution lemma gives
\[
 \Gamma_\sigma\vdash^c g[u:=b]:
       (\uses{t})\theta\le_c(A^\uparrow\cup\{u\})\theta.
\]
This is now evidence over the store alone. Its type still mentions the
substitution of the old use set, which is not yet the displayed use set
of the new body.

\textbf{3. Identify both endpoints.}
Two syntactic identities make the connection:
\[
 \uses{t[u:=b]}=(\uses{t})\theta,\qquad
 (A^\uparrow\cup\{u\})\theta=A\cup\{r\}.
\]
The first is a proved substitution lemma for this particular use-set
function. It is not a consequence of the definition of substitution
alone. For an application, it follows because substitution changes
each operand's root by $\theta$; wrappers add no uses. For a let,
substitution acts on both the right-hand side and its declared allowance
$U$. Value literals have empty immediate uses on both sides.
The second identity removes the parameter binder: $A$ refers only to
outer names, and the singleton parameter becomes $\{r\}$.
Thus the substituted certificate has exactly the endpoints
\[
 \Gamma_\sigma\vdash^c g[u:=b]:
       \uses{t[u:=b]}\le_c A\cup\{r\}.
\]

\textbf{4. Interpret evidence and the stored annotation.}
Capture canonical forms over the typed store turn this judgment into
root inclusion. Separately, store typing gives
$\roots{\{h\}}=\roots{A}$, since $A$ is the annotation of the value
stored at $h$. Combining them gives
\[
 \roots{\uses{t[u:=b]}}
 \subseteq\roots{A\cup\{r\}}
 =\roots{\{h,r\}}
 =\roots{\uses{h\,b}}.
\]
The application step replaces $h\,b$ by $t[u:=b]$, so this is precisely
the required decrease for that step. The inclusion comes from $g$;
the middle equality comes from the stored annotation. Neither is
ordinary subset inclusion between the original sets of program names.

For example, let $b=\cast{\ev{unit}}{\ev{refl}(\ty{Unit})}$ in the
earlier closure. The reduct is $\ev{log}\ b$, not
$\ev{log}\ \ev{unit}$ as target syntax. Nevertheless its use set is
$\{\ev{log},\ev{unit}\}$: the substitution lemma uses
$r=\ev{unit}$. Its roots are $\{\kappa_\ell\}$, exactly as in the
bare-argument trace. This illustrates why the proof keeps the whole
argument in the body and only its root in the capture set.

These are the steps of \texttt{Store.Typed.app\_uses} in the development:
typed substitution, \texttt{Tm.uses\_substAtom}, the closing-set identity,
\texttt{cap\_canon}, and \texttt{Store.Typed.root\_annot}.

\subsection*{Applying the bound through a wrapped receiver}
Now start from a typed application at its \emph{presented} function type:
\[
 \Gamma_\sigma\vdash a:((u:S)\to V)^C,\qquad
 \Gamma_\sigma\vdash b:S.
\]
Let $h=\ev{root}(a)$ and
$\sigma(h)=\lambda^A(u:T).t\langle g\rangle$.
The client's domain $S$ may differ from the stored domain $T$.
After an outer call with argument $b_0$, the generated example has this
form with $a=L$, $b=B[u:=b_0]$, and $h=\ev{log}$.
The use-set clause gives
$\uses{a\,b}=\{h,\ev{root}(b)\}$. We must show that the particular
reduct of the wrapped call has roots below this set.

Canonical forms supply the typing facts needed by the machine's two
wrapped application rules. If the receiver's chain has identity form,
or the capture extension's conversion form preserving the function
shape, the argument already has the stored domain type. The reduct is
$t[u:=b]$, and the preceding four-step argument applies directly.

If the chain instead has function form $\form{pi}(d,c)$, the machine
adapts the argument and result as in Section~\ref{sec:application}.
Write $b'=\cast{b}{d}$ and let $c_b$ be the result evidence instantiated
at the original argument $b$. The reduct is
\[
 t_{\mathrm{next}}=\cast{t[u:=b']}{c_b}.
\]
Typed canonical forms establish $d:S\le T$, hence $b':T$.
The closing-certificate argument therefore applies with $b'$.
The atom-root and term-use clauses then give
\[
 \ev{root}(b')=\ev{root}(b),\qquad
 \uses{t_{\mathrm{next}}}=\uses{t[u:=b']}.
\]
Consequently,
\[
 \roots{\uses{t_{\mathrm{next}}}}
 \subseteq\roots{\{h,\ev{root}(b')\}}
 =\roots{\{h,\ev{root}(b)\}}
 =\roots{\uses{a\,b}}.
\]
The use set records the receiver's root $h$, whose capabilities come
from its stored annotation. The presented capture annotation on $a$'s
type need not be identified with that annotation.

For the generated example's next call, the relevant stored lambda is
$\ev{log}$, so we use its closing certificate. Its identity body has
certificate $\ev{elem}(\{u\},\{\kappa_\ell,u\})$.
The outer certificate $g_{\mathrm{gen}}$ bounded entry into $L\,B$;
the certificate of $\ev{log}$ bounds execution of the subsequently
invoked identity body. Each call uses the evidence belonging to the
closure it actually invokes.

These application steps keep the store and pending continuation $K$
unchanged. A whole state's use set is the union of its current term's
uses with the pending frames' uses. Adding the same $\roots{\uses{K}}$
to both sides therefore gives the application case for a whole state.
The next example shows how those frames acquire their allowances and
why the allowances matter when a continuation is later discharged.

\exercise{roots}{Assume $k$ has capture annotation $\{h,\ev{console}\}$.
Compute $\roots{\{k,\ev{unit}\}}$. Can evidence
$\{k\}\le_c\{\ev{log}\}$ be sound over this typed store when
$\kappa_{\ell}$ and $\kappa_c$ are distinct? Explain using root inclusion.}

\lesson{Immediate uses and result captures}\label{sec:latent}
\focus{A call's immediate uses can differ from the captures of its returned closure.}

To distinguish immediate uses from result captures, consider a factory
that returns a closure associated with another capability.
Figure~\ref{fig:factory} presents the factory and two clients in
simplified target notation, including their closing and avoidance
certificates. Calling $\ev{make}$ invokes $\ev{log}$ and constructs the
inner closure; its call to $\ev{console}$ remains latent until the
returned closure is invoked.


\begin{figure}[!hbp]
\begin{examplepanel}{tgtcol}{Factory and certificates in simplified target notation}
\small
\[
\tgt{\begin{aligned}
 v_{\mathrm{in}}
   &=\lambda^{\{\ev{console}\}}(v:\ty{Unit}).\
      \ev{console}\ v\ \langle g_{\mathrm{in}}\rangle,\\
 v_{\mathrm{make}}
   &=\lambda^{\{\ev{log}\}}(u:\ty{Unit}).\\[-2pt]
   &\quad
      \bigl(\ev{let}\ r=\ev{log}\ u\ \ev{in}\
         v_{\mathrm{in}}\ \langle\varnothing;q_0\rangle\bigr)
       \langle g_{\mathrm{make}}\rangle .
\end{aligned}}
\]
The symbol \(v_{\mathrm{in}}\) abbreviates the lambda literal itself,
with empty immediate uses. Its certificates and those of the factory are:
\[
\tgt{\begin{aligned}
 g_{\mathrm{in}}&=\ev{elem}(\{\ev{console},v\},\{\ev{console},v\}),\\
 q_0&=\ev{elem}(\varnothing,\varnothing),\\
 g_{\mathrm{make}}&=\ev{elem}(\{\ev{log},u\},\{\ev{log},u\}).
\end{aligned}}
\]
In a let annotation \(\langle U;q\rangle\), \(U\) bounds the continuation
body's uses and \(q\) proves the avoidance obligation.
The named certificates are checked under their displayed binders:
\(g_{\mathrm{in}}\) under \(v\), \(g_{\mathrm{make}}\) under \(u\),
and \(q_0\) under \(r\).
\end{examplepanel}
\begin{examplepanel}{tgtcol}{Clients of the stored factory \(\ev{make}\)}
\small
\[
\tgt{\begin{aligned}
 P_{\mathrm{call}}
   &=\ev{let}\ f=\ev{make}\ \ev{unit}\ \ev{in}\
      f\ \ev{unit}\ \langle U;q_{\mathrm{call}}\rangle,\\
 P_{\mathrm{drop}}
   &=\ev{let}\ f=\ev{make}\ \ev{unit}\ \ev{in}\
      \ev{unit}\ \langle\varnothing;q_{\mathrm{drop}}\rangle,\\
 U&=\{\ev{console},\ev{log}\},\\
 q_{\mathrm{call}}&:\{f,\ev{unit}\}\le_c U,\qquad
 q_{\mathrm{drop}}=\ev{capvar}(\ev{unit}):\{\ev{unit}\}\le_c\varnothing.
\end{aligned}}
\]
\end{examplepanel}
\begin{examplepanel}{runcol}{Initial roots of the whole state}
\[
\run{P_{\mathrm{call}}:\{\kappa_\ell,\kappa_c\},
      \qquad P_{\mathrm{drop}}:\{\kappa_\ell\}.}
\]
\end{examplepanel}
\caption{A factory with distinct immediate uses and result captures.
Both clients call \(\mathsf{make}\), which invokes \(\mathsf{log}\) and
returns a closure that invokes \(\mathsf{console}\);
\(P_{\mathrm{call}}\) then invokes that result.
The text constructs \(q_{\mathrm{call}}\) from two branches. This named
presentation follows the C1/C6 examples, with certificates supplied
directly as target evidence. Source translation may produce different
evidence syntax.}
\label{fig:factory}
\end{figure}



The outer closure can therefore have capture annotation $\{\ev{log}\}$,
while its result type has annotation $\{\ev{console}\}$:
\[
 \ev{make}:
 \bigl(\ty{Unit}\to
       (\ty{Unit}\to\ty{Unit})^{\{\ev{console}\}}\bigr)^{\{\ev{log}\}}.
\]
This typing relies on the empty immediate use set of a value expression,
even when its capture annotation is nonempty. The returned lambda's body
is checked with its own closing evidence and executes only when the
closure is called.

\remember{Capture annotations, immediate use sets, and syntactic free
variables are distinct. Effect safety bounds the capability roots
of variables inspected during execution; references reachable from a
returned closure or elsewhere in the heap require a separate property.}

\FloatBarrier
\subsection*{Continuation allowances}
Both clients in Figure~\ref{fig:factory} call $\ev{make}$, then differ in
how their let bodies use its result. The calling client invokes the
returned closure; the dropping client returns $\ev{unit}$. A bound based
only on $\ev{make}$'s annotation would omit the later console call in
$P_{\mathrm{call}}$, while charging every result capture to the initial
call would also charge $P_{\mathrm{drop}}$ for that unused capability.

The continuation supplies the additional information. While evaluating
$\ev{make}\ \ev{unit}$, the machine retains a frame for the pending
body. This frame declares an outer use set $U$ and carries evidence
that bounds the body's uses by $U$.

\needspace{12\baselineskip}
For $P_{\mathrm{call}}$, choose $U=\{\ev{console},\ev{log}\}$.
The continuation binds $f$ with annotation $\{\ev{console}\}$, allowing
the avoidance evidence to be constructed in two branches:
\[
 \begin{aligned}
 \{f\}&\xrightarrow{\ev{capvar}(f)}\{\ev{console}\}
          \xrightarrow{\ev{elem}}U,\\
 \{\ev{unit}\}&\xrightarrow{\ev{capvar}(\ev{unit})}\varnothing
          \xrightarrow{\ev{elem}}U.
 \end{aligned}
\]
Writing $q_f$ and $q_u$ for the two branch expressions, the certificate
$q_{\mathrm{call}}=\ev{union}(q_f,q_u)$ proves
$\{f,\ev{unit}\}\le_c U$. Its outer endpoint $U$ avoids the local name
$f$ by using names that account for its captures.

The chosen overapproximation $U$ retains $\ev{log}$, making the decrease
visible when the calling frame is discharged. The body of
$P_{\mathrm{drop}}$ uses only $\ev{unit}$, so
$U=\varnothing$ suffices, with evidence
$q_{\mathrm{drop}}=\ev{capvar}(\ev{unit}):
\{\ev{unit}\}\le_c\varnothing$.

Inside the factory, the let body is the inner lambda literal. Its empty
use set requires only
$q_0=\ev{elem}(\varnothing,\varnothing)$.
The factory body's use set is consequently
$\{\ev{log},u\}\cup\varnothing=\{\ev{log},u\}$.
This equals the allowance justified by the closing certificate
$g_{\mathrm{make}}$.

\subsection*{Whole-state roots along execution}
Before the factory call, the current expression contributes
$\{\ev{make},\ev{unit}\}$. The pending body contributes its declared $U$.
The initial roots are therefore
\[
 \begin{aligned}
 P_{\mathrm{call}}&:\ \{\kappa_{\ell},\kappa_c\},\\
 P_{\mathrm{drop}}&:\ \{\kappa_{\ell}\}.
 \end{aligned}
\]
Figure~\ref{fig:factory-run} follows the calling client through execution.


\begin{figure}[!htbp]
\begin{examplepanel}{runcol}{Whole-state roots for the calling client}
\small
\renewcommand{\arraystretch}{1.3}
\begin{tabularx}{\linewidth}{@{}>{\raggedright\arraybackslash}Xll@{}}
\toprule
Stage & Whole-state roots & Inspection\\
\midrule
Call \(\mathsf{make}\); keep the calling frame
 & $\{\kappa_\ell,\kappa_c\}$ & $\mathsf{make}$\\
Run its call to \(\mathsf{log}\); keep that frame
 & $\{\kappa_\ell,\kappa_c\}$ & $\mathsf{log}$\\
Inner closure ready for allocation
 & $\{\kappa_\ell,\kappa_c\}$ & none\\
Call the newly allocated \(f\); the frame is discharged
 & $\{\kappa_c\}$ & $f$\\
Run \(\mathsf{console}\ \mathsf{unit}\)
 & $\{\kappa_c\}$ & $\mathsf{console}$\\
Return \(\mathsf{unit}\)
 & $\varnothing$ & none\\
\bottomrule
\end{tabularx}
\end{examplepanel}
\caption{Execution of \(P_{\mathrm{call}}\) from
Figure~\ref{fig:factory}, with administrative transitions grouped into
rows. Before allocation, the pending frame contributes
\(U=\{\mathsf{console},\mathsf{log}\}\).
After that frame is discharged, the call's use set resolves to the console
root alone. The dropping client \(P_{\mathrm{drop}}\) allocates the result
and returns \(\mathsf{unit}\) without invoking it.}
\label{fig:factory-run}
\end{figure}



The trace groups administrative transitions. Discharging the calling
frame replaces its overapproximation, which retains $\ev{log}$, with
the use set of the actual call.

The dropping client also produces and allocates the inner closure, then
returns $\ev{unit}$ without invoking it. The console root is absent from
the initial root set and never occurs among the roots of an inspected
variable.

\exercise{latent}{Suppose the calling client instead tried to declare
$U=\{\ev{log}\}$. Which branch of its avoidance proof would fail?
Why does the same objection not apply to the dropping client?}

\lesson{Capture prediction and its scope}\label{sec:capture-scope}
\focus{Effect safety concerns runtime variable inspection; scope escape and classification require further hypotheses.}

Capture prediction states that the roots of a typed state's use set do
not grow along execution, with renaming accounting for store extension
by allocation.
A second lemma establishes that the roots of each inspected variable
are included in the roots of the current use set. Together these imply
effect safety: a capability root absent from the roots of the initial
use set cannot be a root of any later inspected variable.

The observable events are variable inspections, such as reading the
function receiver of an application or the object receiver of a
projection. Relating these observations to operating-system effects
would require semantics for those effects. Unboxing, which opens a
package that hides captures, has a separate corresponding safety result.

\subsection*{Capture members and inconsistent bounds}
Capture members expose lower and upper bounds, paralleling type members.
The name $x.C$ selects a capture member, as $x.A$ selects a type member;
$\{x.C\}$ is a capture set containing that selection.
Under a hypothetical interface
\[
 x:\{C:\{\kappa_{\ell}\}\,..\,\varnothing\},
\]
reading and composing its bounds gives
\[
 \{\kappa_{\ell}\}\le_c\{x.C\}\le_c\varnothing.
\]
This evidence is derivable under the hypothetical assumption. If a
typed store realized the interface with rigid
$\kappa_{\ell}$, its semantic interpretation would imply
$\{\kappa_{\ell}\}\subseteq\varnothing$, a contradiction.
The argument follows the same pattern as inconsistent type bounds:
the derivation is conditional on an interface that no typed store can
realize.

\needspace{8\baselineskip}
\subsection*{Scope escape}
An effect bound constrains what a computation may inspect. Scope rules
add conditions on whether a locally introduced capability can be returned
or subsumed into an outer capability set. Establishing such conditions
requires an additional argument beyond effect safety.

The \texttt{CapturesCC} extension introduces scope roots, levels, and
fresh capabilities. Its no-escape results depend on explicit restrictions;
in particular, relevant source results restrict derivations that read
capture bounds. These results exclude particular callback subcapturing
judgments under stated assumptions about scopes and bound reads.
Extending them to general \texttt{withFile} programs would require a
corresponding argument for those programs.

\subsection*{Classified capture bounds}
The \texttt{Classifiers} extension labels capabilities with classifiers,
such as a class of control capabilities. A kind specifies the permitted
classifiers, and kinding evidence establishes that a capture set belongs
to that kind. Filtered capture sets describe the corresponding roots;
classified effect safety ensures that inspected roots satisfy the
permitted classification.

This extends the root-set argument with kinding evidence and its
interpretation while retaining the store, substitution, and prediction
reasoning. Operational results for exception handlers or thread
scheduling would additionally require models of those mechanisms.

\lesson{Extensions to the object calculus}\label{sec:larger}
\focus{Fields, recursive types, and intersections introduce additional evidence and normalization cases.}

\subsection*{Field presence and field types}
An object can hold a field $a$ as well as define a type $A$.
A bound on the field's type does not by itself prove that the field
exists. The target also associates type names with field labels:
$s.a$ names the type of field $a$'s body. This is a \emph{block name},
like $s.A$ for a type member; here it denotes a type, not the term that
projects the field's value.
The target therefore records two separate propositions:
a presence fact for $a$, and an inclusion from $s.a$ to its advertised
type.

A literal's precise information includes the field's actual type
definition and its presence. Presence evidence lets the safety proof
establish that field lookup succeeds. Projection runs the stored field
body, and inclusion evidence presents its result at the interface's
declared type.
As with type members, the interface records propositions justified by
the stored literal.

For example, consider an object defining $A=W$ and holding
$a=f$. A client may see its field at type $x.A$ even if the stored field
was originally checked at $W$. The precise facts $s.a\equiv W$ and
$s.A\equiv W$ connect the field's type with the type member.
At $x$, these facts justify presenting the field result at $x.A$.
Field projection and the coercion of its result are separate operations,
each with explicit evidence.

\subsection*{Recursive self types}
A source self type can mention its own members:
an object might advertise that its field $a$ has type $s.A$.
Opening this interface at $x$ replaces the bound self by $x$.
Closing it again restores the self binder.
The target uses fold and unfold wrappers on atoms to record these
changes of presentation. Like casts, the wrappers preserve the root.

The self binder in a telescope allows propositions to refer to their
object. Recursive source typing additionally requires rules for opening
and closing the interface at a particular root.

\subsection*{Intersections and self-bounds}
An intersection $S\wedge T$ describes a value usable at both $S$ and $T$.
The field interface $\{a:T\}$ promises a field $a$ with type $T$.
When both components list member declarations, their intersection can
be represented by collecting facts from both lists.
But the source also allows $x.A\wedge\{a:T\}$, where the selection
$x.A$ supplies no visible list of declarations.

To handle components without a visible telescope, the translation
introduces a \emph{self-bound}: a proposition
saying that the whole self has some type. In this example, the $x.A$
component contributes the fact that the whole self can be used at $x.A$;
the field component contributes presence and a type bound for $a$.

Self-bounds require additional head forms and entries that record
conversions through a bound to the interface used by an operation.
These cases extend the five forms in Section~\ref{sec:normalization};
their typing still ensures that normalized evidence supplies the
information required by the operation.

\subsection*{Cyclic aliases}
An object's type definitions may say $A=s.B$ and $B=s.A$.
Unbounded traversal of these definitions would diverge. The resolver
uses a finite bound derived from the available definitions and resolves
a cycle to $\top$, which contributes no additional information.
Resolving it to $\bot$ would instead assert an uninhabited value type.
This result describes the resolved head shape; it supplies neither a
source typing rule nor equality evidence between a cyclic name and
$\top$.

The bounds for alias resolution and evidence normalization have distinct
justifications. Alias resolution uses finiteness of the available names;
evidence normalization uses induction on typing to establish the
existence of sufficient fuel.

\subsection*{Related developments}
The repository also contains a paths variant, where a field path such as
$x.a.b$ can be the receiver of a type selection $(x.a.b).A$.
Singleton types there record that a value is the one named by a
particular path.
A separate OOPSLA 2016 development has recursive subtyping under a self
assumption and its own evidence target.
The theorems apply to these developments separately; they do not cover
their unrestricted combination. Such a combination would need to
establish field presence, stable path resolution, and recursive
subtyping evidence under the appropriate self assumptions within one
evidence language and machine semantics.

\lesson{The Lean development}\label{sec:lean}
\focus{The companion instantiates the formal syntax and checks the worked calculations.}

The companion \texttt{paper/TutorialExamples.lean} formalizes the core
example using the development's \FC{} definitions. It checks the interface
coercions, store typing, views, member normalization, composition,
and function coercion.
It also constructs DOT-MNF derivations: DOT in \emph{monadic
normal form}, where intermediate computations are named by lets.
These examples check the focused translations from Section~\ref{sec:evidence}
and the complete program from Section~\ref{sec:closed-program}, including
its literal, recursive opening, lambda, and four lets.
From the repository root, run:
\begin{lstlisting}
lake env lean lean/Coercions/paper/TutorialExamples.lean
\end{lstlisting}
The imported libraries must already have been built. In a fresh checkout,
run \texttt{lake build} first.
The equalities are theorems checked by Lean's proof-checking kernel.

\subsection*{From recipe notation to constructors}
The declaration \texttt{def name : Type := expression} defines a named
expression of the stated Lean type.
Here \texttt{LeCo s} is the syntax of inclusion evidence whose variables
belong to scope \texttt{s}. This scope parameter is distinct from the
mathematical self name $s$ used in our telescopes.
The Lean type describes which expressions can be built; the evidence
typing judgment separately checks their inclusion endpoints.

Here is the companion's definition of the exposing coercion:
\begin{lstlisting}
def expose : LeCo s :=
  .obj litTel
    (.le
      (.le .nil .none (.eqSym 0) .none)
      .none (.eq 0) .none)
\end{lstlisting}
The leading dot lets Lean infer the constructor's namespace.
The \texttt{.obj} constructor combines the input telescope
\texttt{litTel} with a morphism. Reading from the inside out,
\texttt{.nil} is the empty morphism and each \texttt{.le} appends an
inclusion recipe. The inner \texttt{.le} supplies the lower recipe,
reading equality slot 0 in reverse with both sides absent.
The outer \texttt{.le} supplies the upper recipe, reading the same slot
forwards.

This is the two-row recipe table from Section~\ref{sec:recipes},
represented as an inductive list. Recovering the list order makes the
individual recipes easy to identify in the constructor expression.

Type selections are represented by scope-indexed variables and labels.
The declaration \texttt{A : Label := .typ 0} chooses a type-member label.
Within the self binder, \texttt{.here} refers to that self, so
\texttt{.sel .here A} represents $s.A$.

\subsection*{What exactly is checked?}
The companion first constructs a store containing a literal defining
$A=W$ and proves its typing. In that store, it establishes the view
equations
\[
 \view{a}=[\ev{le}\,\form{id},\ev{le}\,\form{id}],
 \qquad
 \view{b}=[\ev{le}\,\form{bot},\ev{le}\,\form{top}].
\]
The normalizer takes an explicit fuel argument; a successful result is
represented by \texttt{some} applied to the computed form.
Each example supplies a concrete fuel bound.
\texttt{CanonicalForms.lean} proves that sufficient fuel exists for
every well-typed input closed over a typed store.

The \texttt{SourceTranslation} namespace applies the translation function
to explicit source derivations. Its \texttt{nestedTranslates} and
\texttt{composedTranslates} theorems establish the nested and composed
casts; further theorems prove target typing, distinct target syntax, and
equal erasure. Those focused derivations use the parameter context of the
translation figures. The \texttt{FullSource} namespace then constructs
the closed source derivation. Its \texttt{programTranslates} theorem
identifies the complete target constructor expression,
\texttt{programTyped} derives its typing, and \texttt{programErases}
checks its erasure. \texttt{openedRoot} and \texttt{openedView} check
the two properties of $R$: preservation of the root and preservation of
the view with two additional units of fuel. The store examples check
the underlying view and normalization calculations. We do not claim a
separately mechanized multistep execution theorem for the whole tutorial
program.

The generated capture certificate is a paper calculation of the
translation clauses in \texttt{Terms.lean} and \texttt{Evidence.lean}.
It retains the generated constructor syntax. It has not been added as a
separate companion theorem.
The operational capture examples instead use hand-written target evidence.
The $h\ \ev{unit}$ trace specializes the application and prediction rules
on paper.

\needspace{5\baselineskip}
The factory follows the mechanized C1/C6 examples in
\texttt{Captures/FCdot/Examples.lean}.
It uses names for their encoded variables. Its trace groups administrative transitions; a row can
represent several machine steps.

\subsection*{Code navigation}
\begin{center}
\small
\renewcommand{\arraystretch}{1.3}
\begin{tabularx}{\linewidth}{@{}>{\raggedright\arraybackslash}p{0.35\linewidth}>{\raggedright\arraybackslash}X@{}}
\toprule
Question & Starting point, under \texttt{lean/Coercions/}\\
\midrule
What are slots and recipes? &
\leanfile{FCdot/Syntax.lean}{FCdot/Syntax.lean}; \texttt{Hole}, \texttt{Morphism}\\
Why are sides self-free? &
\leanfile{FCdot/Typing.lean}{FCdot/Typing.lean}; morphism typing\\
How are recipes combined? &
\leanfile{FCdot/Normalizer.lean}{FCdot/Normalizer.lean}; \texttt{Form.combine}, \texttt{LocalEntry.at}\\
Why do the calculations succeed? &
\leanfile{FCdot/CanonicalForms.lean}{FCdot/CanonicalForms.lean}; \texttt{le\_canon}, \texttt{atom\_canon}\\
How is a cast function called? &
\leanfile{FCdot/Machine.lean}{FCdot/Machine.lean}; application steps\\
Does the evidence checker agree? &
\leanfile{FCdot/CheckerCompleteness.lean}{FCdot/CheckerCompleteness.lean}\\
How is source safety obtained? &
\leanfile{DotToFCdot/Safety.lean}{DotToFCdot/Safety.lean}; \texttt{dot\_safety}\\
How is closing evidence generated? &
\leanfile{Captures/DotToFCdot/Terms.lean}{Captures/DotToFCdot/Terms.lean};
\texttt{translate}, \texttt{translateUses}\\
What do captures mean operationally? &
\leanfile{Captures/FCdot/Prediction.lean}{Captures/FCdot/Prediction.lean};
\texttt{app\_uses}, \texttt{capture\_prediction}, \texttt{effect\_safety}\\
Where is the factory example? &
\leanfile{Captures/FCdot/Examples.lean}{Captures/FCdot/Examples.lean};
\texttt{C1\_typed}, \texttt{C6\_run}, \texttt{C6\_safe}\\
\bottomrule
\end{tabularx}
\end{center}
\normalsize

The reference report \emph{Explicit Evidence for Dependent Object
Types} gives the full syntax and core metatheory, including the target
typing rules and canonical-forms theorem used here. The capture report develops use
sets, scopes, and classifiers. Together, the reports and the table above
connect each worked construction to its general rule and Lean declaration.

\exercise{lean}{In the displayed definition of \texttt{expose}, replace
the first \texttt{eqSym 0} by \texttt{eq 0}. What two facts would the
modified recipes produce? Why would that no longer check as an
inclusion from $\ty{Lit}$ to $\ty{Exact}$?}

\lesson{Connecting the model to a compiler}\label{sec:compiler}
\focus{Explicit evidence separates source type reconstruction from validation of its results.}

A source type checker reconstructs a typing derivation by choosing rule
instances and intermediate types and inferring missing annotations.
The target checker validates the resulting evidence.
For composition, the evidence supplies both components; checking
computes their endpoints and verifies that the intermediate types agree.

The target checker is proved sound and complete for its evidence
judgment: it accepts exactly the evidence admitted by that judgment.
This completeness result concerns evidence checking. Completeness of
source inference depends on the source frontend's search procedure.
The repository's frontends produce source derivations and compose with
translation and checking; their search behavior is a separate property.

This boundary supports two experiments with a real compiler.
First, for a specified source fragment, compare the compiler's decisions
with those of the modeled frontend on generated programs and regression
tests. Retain successful derivations and translations. Disagreements
can then be classified as differences in language rules, search
limitations, encoding errors, or compiler defects.

Second, instrument the compiler to record how it discharges typing
obligations, then translate those derivations into independently checked
evidence.
The resulting guarantee is relative to the modeled language and the
translation's stated coverage. Establishing end-to-end compiler
correctness would additionally require proofs that the compiler-to-model
translation and generated code preserve the source semantics.

\subsection*{A nominal source calculus}
A nominal source calculus would add class identity, declared
inheritance, and member lookup. Its elaboration would justify these
declarations in the structural target or in an extension with explicit
nominal evidence.

An inheritance clause generates obligations for inherited members;
a type alias generates an equality obligation; a method override
generates argument and result adaptation obligations.
These extend constructions already encountered in the tutorial.
Nominal identity, initialization, dispatch, and recursive class
structure introduce further semantic requirements.

Such an elaboration must preserve typing and operational behavior,
while retaining the evidence normalization and substitution properties
used by the safety proof.
The existing path, capture, and recursive-subtyping developments provide
components for this investigation. Their integration into a realistic
Scala model remains an open development task.

\remember{The development provides an evidence language, verified
translations for specified source calculi, and safety proofs.
A compiler connection must identify its source fragment, translate its
typing obligations, and state the guarantee obtained from checking them.}

\clearpage\lesson{Worked solutions}\label{sec:answers}
\focus{The key checks are endpoint agreement, root inclusion, and the hypotheses of the normalization and safety theorems.}

\answer{bounds}
The client needs $W<:x.A$. Its available lower bound is $\bot<:x.A$,
so using that bound would also require $W<:\bot$.
The abstract interface provides no such derivation. Although the stored
definition is $A=W$, this interface hides its connection to $W$.

\answer{order}
The reverse composition is
\[
 x.A\xrightarrow{h_x}W\xrightarrow{\ell_x}x.A,
\]
so it proves $x.A\le x.A$, whereas the original composition proves
$W\le W$. The expressions therefore have different endpoint pairs.
Evidence for $W\le\top$ through the member is
$\ell_x;h_x;\ev{top}(W)$.
Alternatively, $\ell_x;\ev{top}(x.A)$ uses the inclusion into $\top$
directly at $x.A$.

\answer{translate}
Let $t_A=\ev{top}(x.A):x.A\le\top$.
Two applications of source subsumption translate to
$\cast{\cast{f}{\ell_x}}{t_A}$.
A source transitivity derivation followed by one subsumption translates to
$\cast{f}{\ell_x;t_A}$.
Each \textsc{Sub} contributes a cast; \textsc{Trans} contributes evidence
composition. Both erase to $f$, although they are distinct constructor
expressions in the target.

\answer{root}
The selected propositions are $\bot\le x.A$ and $x.A\le\top$.
The casts determine the available interface; their root remains $x$,
so both propositions refer to $x.A$.
Removing the casts preserves the root but may invalidate access to
these telescope slots.

\answer{widen}
Use
\[
 \begin{aligned}
 \text{lower: }&
 [\,\ev{bot}(W)\mid\text{inclusion }0\mid\text{absent}\,],\\
 \text{upper: }&
 [\,\text{absent}\mid\text{inclusion }1\mid\text{absent}\,].
 \end{aligned}
\]
The lower chain is $\bot\le W\le s.A$; the upper chain is $s.A\le W$.
Relative to $\ev{forget}$, the upper recipe omits its right side
$\ev{top}(W)$, retaining the stronger upper bound.

\answer{views}
Through $a:\ty{Exact}$, the upper inclusion is $\ell.A\le W$ and its head
form is $\form{id}$. Through $b:\ty{Abs}$, the upper inclusion is
$\ell.A\le\top$ and its head form is $\form{top}$.
The inclusions have different target endpoints while referring to the
same member of the stored object.

\answer{compose}
Interpreting the reversed equality hole against the equality marker
yields $\form{id}$. The left side normalizes to $\form{bot}$, and the
absent right side contributes $\form{id}$. Thus the result is
$\form{bot};\form{id};\form{id}=\form{bot}$.
The unfilled recipe remains parameterized by a source proposition;
supplying the stored object's view instantiates it.

\answer{variance}
The inclusion is
\[
 (S_1\to T_1)\le(S_2\to T_2).
\]
The client supplies $b:S_2$. Evidence $d$ casts the argument to
$\cast{b}{d}:S_1$ for the original body, whose result has type $T_1$.
Evidence $c$ casts that result to $T_2$ for the client.
The domain is contravariant because the argument supplied at the target
domain $S_2$ must be accepted at the source domain $S_1$.

\answer{bad}
The canonical-forms theorem requires evidence typed in the context of
a \emph{typed store}. The open derivation instead assumes $y:\ty{Bad}$
and establishes $\top\le\bot$ conditionally on that assumption.
Consistency over typed stores excludes any stored argument that could
discharge it. The missing hypothesis is store typing, which is stronger
than context well-formedness.

\answer{proof}
For (a), target safety and operational correspondence transfer safety
to the source. Preservation of typing by translation supplies only the
static part of this argument.
For (b), preservation maintains the typing invariant after each step,
allowing progress to apply at every subsequent state.

\answer{roots}
Following $k$'s annotation $\{h,\ev{console}\}$ resolves $h$ through
$\ev{log}$ to $\kappa_{\ell}$ and $\ev{console}$ to $\kappa_c$.
Since $\ev{unit}$ contributes no roots,
\[
 \roots{\{k,\ev{unit}\}}=\{\kappa_{\ell},\kappa_c\}.
\]
By soundness over the typed store, evidence
$\{k\}\le_c\{\ev{log}\}$ would require
$\{\kappa_{\ell},\kappa_c\}\subseteq\{\kappa_{\ell}\}$.
This contradicts the assumed distinctness of the roots.

\answer{latent}
The branch for $f$ fails. Its capture annotation yields
$\{f\}\le_c\{\ev{console}\}$, but the next step to
$\{\ev{log}\}$ would require $\kappa_c$ to belong to
$\{\kappa_{\ell}\}$. The dropping client's body does not use $f$,
so its avoidance proof has no such branch.
The use of $\ev{console}$ remains latent in the constructed closure.

\answer{lean}
Both recipes would read equality slot 0 forwards, producing
$s.A\le W$ twice.
The requested lower proposition of $\ty{Exact}$ is $W\le s.A$.
The modified expression instead describes an interface with a duplicated
upper bound. The checker compares the generated and requested target
telescopes, so this mismatch causes it to reject the proposed inclusion
from $\ty{Lit}$ to $\ty{Exact}$.

\bigskip
\textbf{Independent reconstruction.}
For the core example, reconstruct the chain leading to the call through
the recovered function:
\[
 \begin{aligned}
 &\text{source derivation}\ \longrightarrow\ \text{complete target}\\
 &\quad\longrightarrow\ \text{allocation and closing substitution}
   \ \longrightarrow\ \text{concrete views}\\
 &\quad\longrightarrow\ \text{enabled application step}.
 \end{aligned}
\]
Include the source object rule and recursive opening, rather than starting
with an assumed exact-interface parameter. Distinguish the source binder
$x$, the allocated location $\ell$, and the presented atoms $a,R(a)$.

Next, assume an atom at $\ty{Bad}$ over a typed store. Construct the
contradictory closed evidence and identify the precise theorem hypothesis
that was unavailable for an opaque parameter at the same type.

For captures, reconstruct the argument that bounds a call's reduct:
\[
 \begin{aligned}
 &\text{source body derivation}\ \longrightarrow\ \text{body and certificate}\\
 &\quad\longrightarrow\ \text{typed argument substitution}
   \ \longrightarrow\ \text{rewritten use-set endpoints}\\
 &\quad\longrightarrow\ \text{root inclusion for the reduct}.
 \end{aligned}
\]
Use a wrapped argument. Identify separately the equality about syntactic
uses, the inclusion obtained from capture evidence, and the equality
obtained from the stored closure's annotation. Then explain the extra
premise needed to obtain effect safety from prediction.

For every arrow, identify a supporting passage and its typing premises.
A result guessed from erasure or familiar lambda-calculus behavior does
not reconstruct the target argument if the required evidence is missing.

\end{document}
